% આ ભાગનો પૂર્ણ સંપાદનીય પાઠ છે. શૈલી, ફોન્ટ, ચિત્રો અને પ્રકરણસંદર્ભ માટે સંપૂર્ણ સ્રોત ZIP તથા reader/full-reader.tex જુઓ. આ એકલી સ્વતંત્ર કમ્પાઇલ થતી મુખ્ય ફાઇલ નથી.

% BEGIN OLP-0601
\olpart{mth}{પદ્ધતિઓ}
\phantomsection\label{guunit:OLP-0601}


\begin{editorial}
  આ ભાગમાં સામાન્ય અને પદ્ધતિસંબંધિત સામગ્રી છે;
  ખાસ કરીને ગણિતના વિદ્યાર્થી ન હોય તેને અજાણી લાગી શકે
  એવી જુદી જુદી સાબિતી-પદ્ધતિઓની સમજૂતીઓ છે. હાલમાં
  તેમાં સાબિતીઓ કેવી રીતે લખવી તે અંગેનું એક પ્રકરણ
  અને ગાણિતિક અનુમાન અંગેનું એક પ્રકરણ છે. આગળ
  તેમનાં વધુ વિભાગો, કસરતો અને ગણિતીય પરિભાષા
  અંગેનું એક પ્રકરણ ઉમેરવાનું પણ આયોજન છે.
\end{editorial}

\begingroup
\makeatletter\def\ol@sectioncs{section}\makeatother

% BEGIN OLP-0602
\olchapter{mth}{prf}{સાબિતીઓ}
\olchapterid{mth}{prf}
\phantomsection\label{guunit:OLP-0602}


\begingroup
\makeatletter\def\ol@sectioncs{section}\makeatother

% BEGIN OLP-0603
\olfileid{mth}{prf}{int}

\olsection{પરિચય}
\phantomsection\label{guunit:OLP-0603}


પ્રારંભિક તર્કશાસ્ત્રના તમારા અનુભવને લીધે તમે કદાચ કોઈ
નિષ્પત્તિ તંત્રથી પરિચિત હશો---સંભવતઃ સ્વાભાવિક નિષ્કર્ષણનું
અથવા ફિચ-શૈલીનું નિષ્પત્તિ તંત્ર, કે પછી સાબિતી-વૃક્ષનું તંત્ર.
આવા તંત્રોમાં તમે કદાચ સાબિતીઓ કરી હશે: સૂત્ર સાબિત કર્યું
હશે અથવા આપેલી દલીલ પ્રામાણ્ય ધરાવે છે તે બતાવ્યું હશે. એ માટે તમે
ઇચ્છિત અંતિમ પરિણામ મળે ત્યાં સુધી તંત્રના નિયમો લાગુ પાડ્યા હતા.
પરંતુ તર્કશાસ્ત્ર \emph{વિશે} વિચારતી વખતે પણ આપણે બાબતો સાબિત કરીએ
છીએ અને મોટા ભાગે નિષ્પત્તિ તંત્ર વાપરતા નથી. હકીકતમાં, આપણે
વિચારીએ છીએ એવી મોટા ભાગની સાબિતીઓ સંપૂર્ણપણે પ્રથમ-ક્રમના તર્કની
ભાષામાં નહીં, પરંતુ અંગ્રેજી જેવી સામાન્ય ભાષામાં (અહીં ગુજરાતીમાં;
ક્યારેક સાંકેતિક ભાષા સાથે) લખાય છે. આવી સાબિતી રચતી વખતે શરૂઆતમાં
મૂંઝવણ થાય: નિષ્પત્તિ તંત્ર વિના કંઈક કેવી રીતે સાબિત કરું?
શરૂઆત ક્યાંથી કરું? મારી સાબિતી સાચી છે તે કેવી રીતે જાણું?

સાબિતી કરવાનો પ્રયાસ કરતાં પહેલાં સાબિતી શું છે અને તે કેવી રીતે
રચાય છે તે જાણવું જરૂરી છે. નામ સૂચવે છે તેમ, \emph{સાબિતી}નો
હેતુ કોઈ બાબત સાચી છે તે બતાવવાનો છે. તેને સંવાદ તરીકે વિચારો:
કોઈ પૂછે કે કોઈ બાબત સાચી છે કે નહીં, દાખલા તરીકે, બે સિવાયની
દરેક અવિભાજ્ય સંખ્યા એકી છે કે નહીં. માત્ર ``હા'' કહેવું પૂરતું
નથી; તે \emph{શા માટે} એવું છે તે જાણવા માગી શકે. આ કિસ્સામાં તમે
તેને સાબિતી આપશો.

રોજિંદી ચર્ચામાં ઉત્તરનો માત્ર સંકેત આપવો અથવા અધૂરો ઉત્તર આપવો
પૂરતો હોઈ શકે. પરંતુ તર્કશાસ્ત્ર અને ગણિતમાં આપણે કડક સાબિતી
ઇચ્છીએ છીએ---કોઈ બાબત \emph{જરાય} શંકા રહે નહીં એ રીતે સાચી છે તે
બતાવવા ઇચ્છીએ છીએ. એટલે આપણી સાબિતીનું દરેક પગલું ન્યાયસંગત હોવું
જોઈએ અને તેનો આધાર યોગ્ય હોવો જોઈએ (એટલે, તમે વાપરો છો તે
પૂર્વધારણા ખરેખર તમે સાબિત કરતા પ્રમેયના વિધાનમાં અપાઈ હોવી જોઈએ,
વાપરેલી વ્યાખ્યાઓ યોગ્ય રીતે લાગુ પાડેલી હોવી જોઈએ, અને જે
આધારો આપો છો તે સાચાં અનુમાનો હોવાં જોઈએ, વગેરે).

સામાન્ય રીતે આપણે કોઈ વિધાન સાબિત કરીએ છીએ. સાબિત કરવાના
વિધાનોને જુદાં નામ આપીએ છીએ: પ્રતિજ્ઞાપન, પ્રમેય, ઉપપ્રમેય અથવા
પરિણામ. પ્રતિજ્ઞાપન સાબિત કરવા યોગ્ય મૂળભૂત વિધાન છે: નોંધવા જેટલું
મહત્ત્વનું, પરંતુ કદાચ ખાસ ઊંડું નહીં કે વારંવાર વપરાતું નહીં.
પ્રમેય મહત્ત્વનું પ્રતિજ્ઞાપન છે. તેની સાબિતી ઘણી વાર કેટલાંક
પગલાંમાં વહેંચાય છે અને ક્યારેક તેને સૌપ્રથમ સાબિત કરનાર વ્યક્તિનું
નામ આપવામાં આવે છે (જેમ કે કૅન્ટરનું પ્રમેય, લેવેન્હાઇમ--સ્કોલેમનું
પ્રમેય) અથવા તે જે બાબત અંગે છે તેનું નામ આપવામાં આવે છે (જેમ કે
પૂર્ણતાનું પ્રમેય). ઉપપ્રમેય એ વધુ મહત્ત્વના પરિણામની સાબિતીમાં
વપરાતું પ્રતિજ્ઞાપન અથવા પ્રમેય છે. ગૂંચવણ એ છે કે ક્યારેક ઉપપ્રમેય
પોતે જ મહત્ત્વનું પરિણામ હોય છે અને તેને રજૂ કરનાર વ્યક્તિના નામથી
પણ ઓળખાય છે (જેમ કે ઝોર્નનું ઉપપ્રમેય). પરિણામ એ બીજા કોઈ પરિણામમાંથી
સહેલાઈથી નીકળતું પરિણામ છે.

સાબિત કરવાના વિધાનમાં ઘણી વાર એવી પૂર્વધારણાઓ હોય છે જે સ્પષ્ટ કરે
છે કે આપણે કઈ જાતની વસ્તુઓ વિશે કંઈક સાબિત કરીએ છીએ. તે ``ધારો કે
$!A$ એ $!B \lif !C$ સ્વરૂપનું સૂત્ર છે'' અથવા ``ધારો કે
$\Gamma \Proves !A$'' એવા શબ્દોથી શરૂ થઈ શકે. આ પ્રતિજ્ઞાપન,
પ્રમેય અથવા ઉપપ્રમેયની \emph{પૂર્વધારણાઓ} છે અને તમારી સાબિતીમાં તમે
તેમને સાચી માની શકો. તે આપણે શું સાબિત કરીએ છીએ તે મર્યાદિત કરે
છે અને જે વસ્તુઓ વિશે વાત કરીએ છીએ તેમનાં નામ પણ રજૂ કરે છે.
ઉદાહરણ તરીકે, તમારું પ્રતિજ્ઞાપન ``ધારો કે $!A$ એ $!B \lif !C$
સ્વરૂપનું સૂત્ર છે'' એમ શરૂ થાય, તો તમે ફક્ત ચોક્કસ જાતનાં
સૂત્રો (એટલે શરતી વિધાનો) વિશે કંઈક સાબિત કરો છો; અને તમારી
સાબિતીમાં આવતું $!B \lif !C$ કોઈ પણ મનસ્વી શરતી વિધાન છે એમ
સમજવાનું છે.
% END OLP-0603
\endgroup


\begingroup
\makeatletter\def\ol@sectioncs{section}\makeatother

% BEGIN OLP-0604
\olfileid{mth}{prf}{str}

\olsection{સાબિતીની શરૂઆત}
\phantomsection\label{guunit:OLP-0604}


પરંતુ શરૂઆત ક્યાંથી કરવી?

તમને કંઈક સાબિત કરવા અપાયું છે, તેથી સાબિતીમાં એ બાબત છેલ્લે
આવવી જોઈએ (અલબત્ત, શરૂઆતમાં તમે તેને સાબિત કરવાની \emph{જાહેરાત}
કરી શકો છો, પરંતુ તેને પૂર્વધારણા તરીકે વાપરવી નહીં). તમે જે
સાબિત કરવા માગો છો તે કાગળના નવા પાનાના તળિયે લખો---આ રીતે
તમારું લક્ષ્ય નજર સામે રહેશે.

પછી, કદાચ તમારી પાસે વાપરી શકાય એવી કેટલીક પૂર્વધારણાઓ હશે
(આગળના વિભાગમાં તમારી સાબિતીના \emph{પ્રકાર}ની વાત કરીશું ત્યારે
આ વધુ સ્પષ્ટ થશે). તેને પાનાની ટોચે લખો અને તે પૂર્વધારણાઓ છે
એવું સ્પષ્ટ જણાવો (એટલે, જો તમે $p$ ધારો છો, તો ``ધારો કે $p$''
અથવા ``માની લો કે $p$'' લખો). છેવટે, પ્રશ્નમાં કેટલીક એવી
વ્યાખ્યાઓ હોઈ શકે જે તમારે જાણવી જરૂરી છે. તમને કોઈ ચોક્કસ
વ્યાખ્યા વાપરવાનું કહેવામાં આવ્યું હોઈ શકે, અથવા પૂર્વધારણાઓમાં
કે તમે મેળવવા માગતા નિષ્કર્ષમાં વિવિધ વ્યાખ્યાઓ આવી શકે.
\emph{તેને લખી લો અને તેમનો અર્થ તમને સમજાયો છે તેની ખાતરી કરો.}

તમે સાબિતી કેવી રીતે ગોઠવો તે પ્રશ્નના સ્વરૂપ પર પણ આધાર રાખશે.
આગળનો વિભાગ વાક્યના પ્રકાર પ્રમાણે સાબિતી કેવી રીતે ગોઠવવી તેની
વિગતો આપે છે.
% END OLP-0604
\endgroup


\begingroup
\makeatletter\def\ol@sectioncs{section}\makeatother

% BEGIN OLP-0605
\olfileid{mod}{prf}{def}

\olsection{વ્યાખ્યાઓનો ઉપયોગ}
\phantomsection\label{guunit:OLP-0605}


અગાઉ આપણે કહ્યું હતું કે સાબિતીમાં વપરાઈ શકે એવી બધી વ્યાખ્યાઓથી
તમારે પરિચિત હોવું જોઈએ અને તેમને યોગ્ય રીતે લાગુ પાડતા આવડવું
જોઈએ. આ ખૂબ અગત્યની વાત છે; તેથી તેને થોડું વધુ વિગતે જોઈએ.
વ્યાખ્યાઓ ગુણધર્મો અને સંબંધોનાં ટૂંકાં નામ આપે છે, જેથી આપણે
તેમના વિશે સંક્ષેપમાં વાત કરી શકીએ. રજૂ થયેલા ટૂંકા નામને
\emph{વ્યાખ્યેય પદ} (definiendum) અને તે જેના માટે વપરાય છે તેને
\emph{વ્યાખ્યાકારક વિધાન} (definiens) કહીએ. સાબિતીમાં આપણે ઘણી વાર
વ્યાખ્યેય પદ કેવી રીતે રજૂ થયું હતું તે તરફ પાછા જવું પડે છે:
સાબિતીને આગળ વધારવા વ્યાખ્યાકારક વિધાનની તાર્કિક રચનાનો ઉપયોગ
કરવો પડે છે. વ્યાખ્યાઓને ઉકેલતાં, તમે તર્કની મુખ્ય કામગીરી ક્યાં
થાય છે ત્યાં પહોંચો છો.

એક ઉદાહરણથી શરૂ કરીએ. ધારો કે તમારે નીચેનું સાબિત કરવું છે:

\begin{prop}
કોઈ પણ ગણો $A$ અને $B$ માટે $A \cup B = B \cup A$.
\end{prop}

સાબિતી શરૂ કરવા માટે પણ બે ગણો એકસમાન હોવાનો અર્થ જાણવો પડે;
એટલે, ઉપરના સમીકરણમાં ગણો માટે ``$=$''નો અર્થ શું છે તે જાણવું પડે.
ગણો ત્યારે જ એકસમાન છે જ્યારે તેમના ઘટકો સરખા હોય.
તેથી આપણે નીચેની વ્યાખ્યા ઉકેલવી પડે છે:

\begin{defn}
ગણો $A$ અને $B$ \emph{એકસમાન} છે, $A = B$, ત્યારે અને ત્યારે જ
જ્યારે~$A$નું દરેક ઘટક~$B$નું ઘટક હોય અને ઊલટું પણ હોય.
\end{defn}

આ વ્યાખ્યામાં $A$ અને~$B$ મનસ્વી ગણો માટેનાં સ્થાનાપન્ન ચિહ્નો છે.
તે જેની વ્યાખ્યા આપે છે---\emph{વ્યાખ્યેય પદ}---તે ``$A = B$'' છે;
અને $A = B$ ક્યારે સાચું છે તેની શરત આપે છે. આ શરત---``~$A$નું દરેક
ઘટક~$B$નું ઘટક છે અને ઊલટું પણ છે''---એ
\emph{વ્યાખ્યાકારક વિધાન} છે.\footnote{આ ખાસ કિસ્સામાં---અને આ
  ગૂંચવણ ઊભી કરે છે!---જ્યારે $A = B$ હોય, ત્યારે $A$ અને $B$
  ખરેખર એક જ ગણ છે, ભલે તેની ડાબી અને જમણી બાજુ માટે જુદા
  અક્ષરો વાપરીએ. પરંતુ તે ગણને દર્શાવવાની રીતો જુદી હોઈ શકે,
  અને તેથી આ વ્યાખ્યા તુચ્છ નથી.} વ્યાખ્યા કહે છે કે શરત પૂરી
થાય ત્યારે અને ત્યારે જ $A = B$ સાચું છે (અંગ્રેજીમાં આને
ટૂંકમાં ``iff'' લખાય છે).

વ્યાખ્યા લાગુ પાડતી વખતે તેમાંના $A$ અને $B$ને તમે જે ખાસ
કિસ્સો તપાસો છો તેની સાથે મેળવો. આપણા કિસ્સામાં,
$A \cup B = B \cup A$ સાચું હોવા માટે દરેક $z \in A \cup B$,
$B \cup A$માં પણ હોવું જોઈએ અને ઊલટું પણ. પ્રતિજ્ઞાપનમાંનું
$A \cup B$ વ્યાખ્યાના~$A$નું સ્થાન લે છે અને $B \cup A$ તેના~$B$નું.
કારણ કે $A$ અને $B$ બંને વ્યાખ્યા અને સાબિત કરવાના પ્રતિજ્ઞાપનમાં
વપરાય છે, પણ તેમની ભૂમિકાઓ જુદી છે, તેથી તેમને ગૂંચવી ન દેવા
ધ્યાન રાખો. દાખલા તરીકે,~$A$નું દરેક ઘટક~$B$નું
ઘટક છે અને ઊલટું પણ છે તે બતાવીને પ્રતિજ્ઞાપન સાબિત
થઈ જશે એમ માનવું ભૂલ હશે: તે $A = B$ બતાવશે, $A \cup B = B \cup A$
નહીં. (વળી, $A$ અને $B$ કોઈ પણ બે ગણો હોઈ શકે; $A$ અને~$B$
વિશે કંઈ ધાર્યું ન હોય તો તેઓ જુદા ગણો પણ હોઈ શકે છે.)

સાબિતીમાં આપણે ગણસિદ્ધાંતના યોગગણ જેવા ખ્યાલો સાથે કામ કરીએ
છીએ; એટલે સાબિતી કેવી રીતે આગળ વધારવી તે સમજવા $\cup$ સંકેતનો
અર્થ પણ જાણવો પડે. ક્યારેક એક વ્યાખ્યા ઉકેલવાથી બીજી વ્યાખ્યાઓ
પણ ઉકેલવી પડે છે. દાખલા તરીકે, $A \cup B$ની વ્યાખ્યા
$\Setabs{z}{z \in A \text{ અથવા } z \in B}$ છે. તેથી $x \in A \cup
B$ સાબિત કરવું હોય, તો $\cup$ની વ્યાખ્યા ઉકેલતાં
$x \in \Setabs{z}{z \in A \text{ અથવા } z \in B}$ સાબિત કરવું પડે.
હવે એ પણ યાદ રાખો કે $x \in \Setabs{z}{\dots z\dots}$ ત્યારે અને
ત્યારે જ જ્યારે $\dots x\dots$. તેથી $\Setabs{z}{\dots z \dots}$
સંકેતની વ્યાખ્યા વધુ ઉકેલતાં તમારે $x \in A$ અથવા $x \in B$
બતાવવું પડે. એટલે ``$A \cup B$નું દરેક ઘટક, $B \cup A$નું
પણ ઘટક છે'' તેનો ખરેખર અર્થ છે: ``દરેક $x$ માટે,
જો $x \in A$ અથવા $x \in B$, તો $x \in B$ અથવા $x \in A$.''
પ્રતિજ્ઞાપનમાંની વ્યાખ્યાઓ પૂરેપૂરી ઉકેલીએ તો આપણે નીચેનું
બતાવવાનું રહે છે:

\begin{prop}
કોઈ પણ ગણો $A$ અને $B$ માટે: (a) દરેક $x$ માટે, જો $x \in A$
અથવા $x \in B$, તો $x \in B$ અથવા $x \in A$; અને (b) દરેક
$x$ માટે, જો $x \in B$ અથવા $x \in A$, તો $x \in A$ અથવા
$x \in B$.
\end{prop}

અગત્યની વાત એ છે કે વ્યાખ્યાઓ ઉકેલવી એ સાબિતી રચવાનો આવશ્યક
ભાગ છે. યોગ્ય રીતે આમ કરવું ક્યારેક મુશ્કેલ હોય છે: વ્યાખ્યામાંના
ચલો અને તમે સાબિત કરતા દાવામાંના પદોને જુદાં ઓળખીને તેમની વચ્ચે
યોગ્ય મેળ બેસાડવો પડે. સફળ થવા માટે પ્રશ્ન શું પૂછે છે અને તેમાં
વપરાયેલા બધા પદોનો અર્થ શું છે તે જાણવું જરૂરી છે---ઘણી વાર
એકથી વધુ વ્યાખ્યાઓ ઉકેલવી પડે છે. નીચેનાં સરળ ઉદાહરણોમાં
ઉકેલ લગભગ તરત જ વ્યાખ્યાઓમાંથી નીકળે છે. અલબત્ત, હંમેશાં
આટલું સહેલું નહીં હોય.

\begin{prob}
ધારો કે તમારે $A \cap B \neq \emptyset$ સાબિત કરવું છે. અહીં
આવતી બધી વ્યાખ્યાઓ ઉકેલો, એટલે કે ``$\cap$'', ``='', કે
``$\emptyset$''નો ઉલ્લેખ કર્યા વિના આ વાત ફરીથી કહો.
\end{prob}
% END OLP-0605
\endgroup


\begingroup
\makeatletter\def\ol@sectioncs{section}\makeatother

% BEGIN OLP-0606
\olfileid{mth}{prf}{pat}

\olsection{અનુમાનના નમૂનાઓ}
\phantomsection\label{guunit:OLP-0606}


સાબિતીઓ વ્યક્તિગત અનુમાનોથી બને છે. આપણે અનુમાન કરીએ ત્યારે
સામાન્ય રીતે ``તેથી'', ``આથી'' કે ``એટલે'' જેવા શબ્દોથી તે
દર્શાવીએ છીએ. અનુમાન ઘણી વાર સાબિતીમાં અગાઉથી ઉપલબ્ધ એક કે બે
હકીકતો પર આધાર રાખે છે---તે કોઈ ધારેલી બાબત હોઈ શકે અથવા
અગાઉના અનુમાનથી મેળવેલો નિષ્કર્ષ હોઈ શકે. સ્પષ્ટતા માટે આપણે
આવી બાબતોને ચિહ્નિત કરી શકીએ અને નવા અનુમાનમાં અન્ય કયાં વિધાનો
વાપરીએ છીએ તે જણાવી શકીએ. ઘણી વાર અનુમાનમાં અગાઉની બાબતોમાંથી
નવો નિષ્કર્ષ \emph{શા માટે} નીકળે છે તેની સમજ પણ હોય છે.
સાબિતીમાં વારંવાર વપરાતા અનુમાનના કેટલાક સામાન્ય નમૂનાઓ નીચે
જોઈશું. ગાણિતિક અનુમાનથી સાબિતી જેવા કેટલાક નમૂનાઓ વધુ વિસ્તૃત
છે (અને તેમની ચર્ચા પછી થશે).

અનુમાનનો એક નમૂનો આપણે જોઈ ચૂક્યા છીએ: વ્યાખ્યા ઉકેલવી અથવા
લાગુ પાડવી. વ્યાખ્યા ઉકેલીએ ત્યારે વ્યાખ્યેય પદ ધરાવતી બાબતને
વ્યાખ્યાકારક વિધાનથી ફરી કહીએ છીએ. દાખલા તરીકે, ધારો કે સાબિતીમાં
આપણે $D = E$ પહેલેથી સ્થાપિત કર્યું છે (a). હવે ગણોની સમાનતા
માટેની $=$ની વ્યાખ્યા લાગુ પાડીને કહી શકીએ: ``આથી, (a) અને
વ્યાખ્યા મુજબ,~$D$નું દરેક ઘટક~$E$નું ઘટક છે
અને ઊલટું પણ છે.''

થોડું ગૂંચવણભર્યું એ છે કે અનુમાન કરતી ઘડીએ તેનો આધાર ઘણી વાર
લખતા નથી, પરંતુ પહેલાં લખીએ છીએ. ધારો કે $D = E$ હજી સાબિત
કર્યું નથી, પણ કરવું છે. જો $D = E$ આપણું લક્ષ્ય હોય, તો
વ્યાખ્યા લાગુ પાડીને લક્ષ્યને ફરી કહી શકીએ: $D = E$ સાબિત
કરવા~$D$નું દરેક ઘટક~$E$નું ઘટક છે અને
ઊલટું પણ છે તે સાબિત કરવું પડે. તેથી સાબિતીનું સ્વરૂપ આવું હશે:
(a)~$D$નું દરેક ઘટક~$E$નું ઘટક છે તે સાબિત
કરો; (b)~$E$નું દરેક ઘટક~$D$નું ઘટક છે તે
સાબિત કરો; (c) તેથી (a), (b) અને $=$ની વ્યાખ્યા પરથી $D = E$.
પરંતુ સામાન્ય રીતે આપણે આવું લખતા નથી. તેના બદલે આવું લખીએ:
\begin{quote}
આપણે $D = E$ બતાવવું છે. $=$ની વ્યાખ્યા મુજબ તેનો અર્થ
એ છે કે~$D$નું દરેક ઘટક~$E$નું ઘટક છે અને
ઊલટું પણ છે તે બતાવવું.

(a) \dots (~$D$નું દરેક ઘટક~$E$નું ઘટક છે
તેની સાબિતી) \dots

(b) \dots (~$E$નું દરેક ઘટક~$D$નું ઘટક છે
તેની સાબિતી) \dots
\end{quote}

\subsection{સંયોજનનો ઉપયોગ}

કદાચ અનુમાનનો સૌથી સરળ નમૂનો સંયોજનના બે ભાગમાંથી એકને
નિષ્કર્ષ તરીકે મેળવવાનો છે. એટલે, જો આપણે $p$ અને~$q$ ધાર્યાં
હોય કે અગાઉ સાબિત કર્યાં હોય, તો~$p$ (અને~$q$ પણ) અનુમાની
શકીએ. આ એટલું મૂળભૂત અનુમાન છે કે તેનો ઘણી વાર ઉલ્લેખ પણ
થતો નથી. દાખલા તરીકે, $D = E$ની વ્યાખ્યા ઉકેલ્યા પછી આપણે
સ્થાપિત કર્યું કે~$D$નું દરેક ઘટક~$E$નું ઘટક
છે અને ઊલટું પણ છે. તેમાંથી~$E$નું દરેક ઘટક~$D$નું
ઘટક છે એવો નિષ્કર્ષ કાઢી શકાય (આ ``ઊલટું પણ'' વાળો ભાગ છે).

\subsection{સંયોજન સાબિત કરવું}

ક્યારેક તમને સાબિત કરવા અપાયેલું વિધાન સંયોજનના સ્વરૂપનું હોય:
``$p$ અને $q$ સાબિત કરો.'' એ વખતે ફક્ત બે કામ કરવાનાં છે:
$p$ સાબિત કરો અને પછી $q$ સાબિત કરો. સ્પષ્ટતા માટે સાબિતીને
બે ભાગમાં વહેંચીને તેમના નામ આપી શકો. શરૂઆતની નોંધમાં
પાનાની ટોચે ``(1) $p$ સાબિત કરો'' અને પાનાની વચ્ચે
``(2) $q$ સાબિત કરો'' લખી શકો. (અલબત્ત, તમને સંયોજન
સાબિત કરવા સીધું ન કહેવાયું હોય, છતાં તમારી સાબિતીમાં તેની
જરૂર પડી શકે. દાખલા તરીકે, $D = E$ સાબિત કરવાનું હોય તો
$=$ની વ્યાખ્યા ઉકેલ્યા પછી તમારે સાબિત કરવાનું થાય છે કે
~$D$નું દરેક ઘટક~$E$નું ઘટક છે
\emph{અને}~$E$નું દરેક ઘટક~$D$નું ઘટક છે.)

\subsection{વિયોજન સાબિત કરવું}

સાબિત કરવાનું વિધાન વિયોજનના સ્વરૂપનું હોય (એટલે ``$p$ અથવા
$q$''ના સ્વરૂપનું વિધાન હોય), તો બે વિકલ્પમાંથી એક સાચો છે તે
બતાવવું પૂરતું છે. પરંતુ તમારા પ્રમેયની પૂર્વધારણામાંથી કોઈ એક
વિકલ્પ સીધો નીકળી આવે એવું લગભગ ક્યારેય બનતું નથી. વધુ વખત
પૂર્વધારણાઓ પોતે વિયોજક હોય છે, અથવા તમે બતાવતા હો છો કે અમુક
જાતની બધી વસ્તુઓમાં બેમાંથી એક ગુણધર્મ હોય છે---પણ કેટલીક
વસ્તુઓમાં પહેલો અને બીજીમાં બીજો ગુણધર્મ હોય છે. આવી સ્થિતિમાં
કિસ્સાવાર સાબિતી ઉપયોગી થાય છે (નીચે જુઓ).

\subsection{શરતી સાબિતી}

તમને મળતા ઘણાં પ્રમેયો શરતી સ્વરૂપના હશે (એટલે, $p$ સાચું હોય
તો $q$ પણ સાચું છે તે બતાવવાનું હોય). આવા કિસ્સાઓમાં સાબિતી
ગોઠવવી સરળ છે---શરતી વિધાનનો પૂર્વપક્ષ (અહીં $p$) ધારો અને
તેમાંથી પરિણામ~$q$ સાબિત કરો. એટલે તમારું પ્રમેય ``જો $p$ તો
$q$'' કહે, તો સાબિતીની શરૂઆત ``$p$ ધારો''થી કરો અને અંતે
~$q$ સાબિત થવું જોઈએ.

શરતી વિધાનો જુદી રીતે કહેવાય છે. ``જો $p$ તો $q$''ને બદલે
પ્રમેય ``$p$ ફક્ત ત્યારે જ જો $q$'', ``જો $p$ તો $q$'' અથવા
``$p$ હોય એ શરતે $q$'' એમ કહી શકે. આ બધાનો અર્થ એક જ છે:
$p$ ધારીને તેમાંથી~$q$ સાબિત કરવું. યાદ રાખો કે દ્વિશરતી
વિધાન (``$p$ ત્યારે અને ત્યારે જ જ્યારે $q$'', અંગ્રેજીમાં ``iff'') ખરેખર બે
શરતી વિધાનો છે: જો $p$ તો $q$, અને જો $q$ તો~$p$. તેથી
શરતી સાબિતીના બે દાખલા પૂરતા છે: પહેલા શરતી વિધાન માટે એક
અને બીજા માટે બીજો. પરંતુ ક્યારેક ``ત્યારે અને ત્યારે જ''નાં
વિધાનોની સાંકળથી પણ દ્વિશરતી વિધાન સાબિત કરી શકાય છે, જેમાં
``$p$''થી શરૂ કરીને ``$q$'' સુધી પહોંચીએ---પરંતુ ત્યારે
દરેક પગલું ખરેખર બંને દિશામાં ચાલે છે તેની ખાતરી કરવી પડે.

\subsection{સાર્વત્રિક દાવા}

સાર્વત્રિક દાવો વાપરવો સરળ છે: કોઈ બાબત બધી વસ્તુઓ માટે
સાચી હોય તો દરેક ચોક્કસ વસ્તુ માટે પણ સાચી છે. દાખલા તરીકે,
સાબિતીની પૂર્વધારણા $A \subseteq B$ હોય, તો $\subseteq$ની
વ્યાખ્યા ઉકેલતાં તેનો અર્થ એ છે કે દરેક $x \in A$ માટે
$x \in B$. તેથી $z \in A$ અગાઉથી જાણતા હો તો
$z \in B$ એવો નિષ્કર્ષ કાઢી શકો.

સાર્વત્રિક દાવો સાબિત કરવો થોડો અઘરો લાગશે. સામાન્ય રીતે આ
વિધાનોનું સ્વરૂપ ``જો $x$માં~$P$ ગુણધર્મ હોય તો તેમાં~$Q$
ગુણધર્મ પણ છે'' અથવા ``બધા $P$ એ $Q$ છે'' એવું હોય છે.
અલબત્ત, સ્વરૂપ હંમેશાં બરાબર એવું ન પણ હોય; તમને ચોક્કસ
શું સાબિત કરવાનું છે તે સમજવામાં થોડી કસરત જોઈએ. પરંતુ ઘણી
વાર કોઈ ગુણધર્મ ધરાવતી બધી વસ્તુઓમાં બીજો ગુણધર્મ પણ છે તે
સાબિત કરવું પડે છે.

સાર્વત્રિક દાવો સાબિત કરવા પ્રથમ ગુણધર્મ ધરાવતી વસ્તુઓને નામ
અથવા ચલ આપો અને પછી બતાવો કે તેમાં બીજો ગુણધર્મ પણ છે.
બીજી રીતે કહીએ તો, \emph{બધા}~$P$ માટે કંઈક સાબિત કરવા
\emph{મનસ્વી}~$P$ માટે તે સાબિત કરો. રજૂ કરાયેલું નામ કોઈ
પણ મનસ્વી~$P$ માટે છે. સામાન્ય રીતે મનસ્વી વસ્તુઓ માટે
એક અક્ષરનાં નામ વાપરીએ છીએ અને અક્ષરોની પરંપરા હોય છે:
જેમ કે પ્રાકૃતિક સંખ્યાઓ માટે $n$, સૂત્રો માટે $!A$,
ગણો માટે $A$, વિધેયો માટે $f$, વગેરે.

મુખ્ય કળા આખી સાબિતીમાં સામાન્યતા જાળવવાની છે. શરૂઆતમાં મનસ્વી
વસ્તુ (``$x$'')માં~$P$ ગુણધર્મ છે એમ ધારો અને ફક્ત વ્યાખ્યાઓ
અથવા માન્ય પૂર્વધારણાઓના આધારે બતાવો કે $x$માં~$Q$ ગુણધર્મ
પણ છે. તેમાં $P$ છે તે સિવાય $x$ ચોક્કસ શું છે તે નક્કી
કર્યું નથી; તેથી $P$ ધરાવતી દરેક વસ્તુમાં~$Q$ છે એમ કહી
શકો. ટૂંકમાં, $x$ એ~$P$ ગુણધર્મ ધરાવતી \emph{બધી} વસ્તુઓનું
પ્રતિનિધિત્વ કરે છે.

\begin{prop}
  બધા ગણો $A$ અને $B$ માટે $A \subseteq A \cup B$.
\end{prop}

\begin{proof}
  $A$ અને $B$ મનસ્વી ગણો લો. આપણે $A
  \subseteq A \cup B$ બતાવવું છે. $\subseteq$ની વ્યાખ્યા મુજબ
  તેનો અર્થ છે: દરેક $x$ માટે, જો $x \in A$ તો $x \in A \cup B$.
  તેથી~$x \in A$ને~$A$નું મનસ્વી ઘટક લો. આપણે $x \in A
  \cup B$ બતાવવું છે. $x \in A$ હોવાથી $x \in A$ અથવા $x \in B$.
  આથી $x \in \Setabs{x}{x \in A \lor x \in B}$. પરંતુ $\cup
  $ની વ્યાખ્યા મુજબ તેનો અર્થ $x \in A \cup B$ થાય છે.
\end{proof}

\subsection{કિસ્સાવાર સાબિતી}

ધારો કે પૂર્વધારણા અથવા અગાઉના નિષ્કર્ષ તરીકે તમારી પાસે વિયોજન
છે---તમે $p$ અથવા $q$ સાચું છે એમ ધાર્યું કે સાબિત કર્યું છે.
તમારે $r$ સાબિત કરવું છે. તે બે પગલાંમાં કરો: પહેલાં $p$ સાચું
છે એમ ધારીને~$r$ સાબિત કરો; પછી $q$ સાચું છે એમ ધારીને ફરી
~$r$ સાબિત કરો. આ રીત ચાલે છે કારણ કે બે વિકલ્પમાંથી એક
સાચો છે એવું આપણે ધારીએ છીએ અથવા જાણીએ છીએ. બંને પગલાં
બતાવે છે કે દરેક વિકલ્પ~$r$ના સત્ય માટે પૂરતો છે. (બંને
વિકલ્પ સાચા હોય તો $r$ સાચું છે તેનાં એક નહીં પણ બે કારણો
મળે છે. $p$ અને~$q$ બંને ધારીને $r$ અલગથી સાબિત કરવાની
જરૂર નથી.) આપણે શું કરીએ છીએ તે જણાવવા ``કિસ્સા અલગ
પાડીએ'' એમ કહીએ છીએ. દાખલા તરીકે, ધારો કે $x \in B \cup C$
છે. $B \cup C$ની વ્યાખ્યા $\Setabs{x}{x \in B \text{ અથવા } x \in C}$
છે. એટલે, વ્યાખ્યા મુજબ $x \in B$ અથવા $x \in C$. તેમાંથી
$x \in A$ સાબિત કરવું હોય, તો પહેલાં $x \in B$ ધારીને આ
પૂર્વધારણાથી $x \in A$ સાબિત કરીએ; પછી $x \in C$ ધારીને
ફરી તેમાંથી $x \in A$ સાબિત કરીએ. પૂર્વધારણા નીચે ``કિસ્સા
અલગ પાડીએ'' લખો; તેની નીચે ``કિસ્સો (1): $x \in B$'' અને
પાનાની વચ્ચે ``કિસ્સો (2): $x \in C$'' લખો. પછી પાનાના
ઉપરના અને નીચેના અર્ધભાગમાં સાબિતીઓ પૂરી કરો.

જે વિધાન સાબિત કરવાનું હોય તે પોતે વિયોજક હોય ત્યારે કિસ્સાવાર
સાબિતી ખાસ ઉપયોગી છે. આ સરળ ઉદાહરણ જુઓ:

\begin{prop}
ધારો કે $B \subseteq D$ અને $C \subseteq E$. તો $B \cup C \subseteq
D \cup E$.
\end{prop}

\begin{proof}
  (a) $B \subseteq D$ અને (b) $C \subseteq E$ ધારો. વ્યાખ્યા
  મુજબ, દરેક $x \in B$ માટે $x \in D$ પણ છે (c), અને દરેક
  $x \in C$ માટે $x \in E$ પણ છે (d).\sourcecorrection{OLMD-155}{મૂળમાં
  (c) અને (d)ના સભ્યતા-વિધાનમાં $x$ છૂટી ગયું છે; અહીં તે
  સ્પષ્ટ કર્યું છે.} $B \cup C \subseteq D \cup E$ બતાવવા
  આપણે બતાવવું છે કે જો $x \in B \cup C$ તો $x \in D \cup E$
  ($\subseteq$ની વ્યાખ્યા મુજબ). $x \in B \cup C$ ત્યારે
  અને ત્યારે જ જ્યારે $x \in B$ અથવા $x \in
  C$ ($\cup$ની વ્યાખ્યા મુજબ). એ જ રીતે, $x \in D \cup E$
  ત્યારે અને ત્યારે જ જ્યારે $x \in D$ અથવા $x \in E$.
  તેથી દરેક $x$ માટે બતાવવાનું છે કે જો $x \in B$ અથવા
  $x \in C$, તો $x \in D$ અથવા $x \in E$.

  \begin{quote}
  અત્યાર સુધી આપણે ફક્ત વ્યાખ્યાઓ ઉકેલી છે!{} આપણે
  પ્રતિજ્ઞાપનને $\subseteq$ અને $\cup$ વિના ફરી કહ્યું છે;
  હવે સાર્વત્રિક શરતી દાવો સાબિત કરવો છે. ઉપરની ચર્ચા મુજબ,
  આ માટે એવું $x$ ધારો જેના વિશે ``જો'' વાળો ભાગ સાચો છે
  અને પછી બતાવો કે ``તો'' વાળો ભાગ પણ સાચો છે. એટલે,
  $x
  \in B$ અથવા $x \in C$ ધારીને $x \in D$ અથવા $x \in
  E$ બતાવીશું.\footnote{આ ફકરો આપણે શું કરીએ છીએ તે જ સમજાવે
    છે---તે સાબિતીનો ભાગ નથી, અને તમે તમારી સાબિતી લખો ત્યારે
    આટલી વિગતમાં જવાની જરૂર નથી.}
  \end{quote}

  ધારો કે $x \in B$ અથવા $x \in C$. આપણે $x \in D$
  અથવા $x \in E$ બતાવવું છે. કિસ્સા અલગ પાડીએ.

  કિસ્સો 1: $x \in B$. (c) પરથી $x \in D$. તેથી $x \in D$ અથવા
  $x \in
  E$. (અહીં અગાઉના ઉપવિભાગમાં ચર્ચાયેલું અનુમાન કર્યું છે!)

  કિસ્સો 2: $x \in C$. (d) પરથી $x \in E$. તેથી $x \in D$ અથવા $x \in E$.
 \end{proof}

\subsection{અસ્તિત્વનો દાવો સાબિત કરવો}

અસ્તિત્વનો દાવો સાબિત કરવા કહેવાય ત્યારે પ્રશ્ન સામાન્ય રીતે
``એવો~$x$ છે કે $\dots x \dots$ તે સાબિત કરો'' એ સ્વરૂપનો
હશે; એટલે ``$\dots x
\dots$''થી વર્ણવાયેલો ગુણધર્મ ધરાવતી કોઈ વસ્તુ છે તે બતાવવાનું
હશે. એ માટે યોગ્ય વસ્તુ શોધીને બતાવવું પડે કે તેમાં જરૂરી
ગુણધર્મ છે. સાંભળવામાં સીધું લાગે છે, પણ આવી સાબિતી મુશ્કેલ
હોઈ શકે. સામાન્ય રીતે વસ્તુનું \emph{નિર્માણ} અથવા તેની
\emph{વ્યાખ્યા} આપવી પડે અને પછી સાબિત કરવું પડે કે એમ
વ્યાખ્યાયિત વસ્તુમાં જરૂરી ગુણધર્મ છે. યોગ્ય વસ્તુ શોધવી
મુશ્કેલ હોઈ શકે; તેમાં જરૂરી ગુણધર્મ છે તે સાબિત કરવું મુશ્કેલ
હોઈ શકે; અને ક્યારેક તમે ખરેખર કોઈ વસ્તુ વ્યાખ્યાયિત કરી છે
એ જ બતાવવું મુશ્કેલ હોય છે!{}

સામાન્ય રીતે તમે વસ્તુ નિર્દેશીને લખો, જેમ કે ``$x$ને
\dots લો'' (અહીં \dots{}થી કઈ વસ્તુ મનમાં છે તે કહો), જરૂર
હોય તો $\dots$ ખરેખર અસ્તિત્વ ધરાવતી વસ્તુ વર્ણવે છે તે
સાબિત કરો, અને પછી બતાવો કે $x$માં~$Q$ ગુણધર્મ છે. એક
સરળ ઉદાહરણ જુઓ.

\begin{prop}
  ધારો કે $x \in B$. તો એવો~$A$ છે કે $A \subseteq
  B$ અને $A \neq \emptyset$.
\end{prop}

\begin{proof}
  $x \in B$ ધારો. $A = \{x\}$ લો.
  \begin{quote}
    અહીં આપણે તેના ઘટકોની યાદી આપીને ગણ~$A$ની
    વ્યાખ્યા કરી છે. $x$ એક વસ્તુ છે એમ ધાર્યું છે અને તેના
    ઘટકોની યાદીથી ગણ હંમેશાં બનાવી શકાય છે; તેથી અહીં
    ગણ~$A$ની સફળતાપૂર્વક વ્યાખ્યા થઈ છે તે અલગથી બતાવવું
    પડતું નથી. પરંતુ પ્રતિજ્ઞાપન માગે છે તે ગુણધર્મો $A$માં
    છે એવું હજી બતાવવું પડે. તેના વિના સાબિતી પૂરી નથી!{}
  \end{quote}
  $x \in A$ હોવાથી $A \neq \emptyset$.
  \begin{quote}
    અહીં $A$ની $\{x\}$ તરીકેની વ્યાખ્યા અને $x \in \{x\}$
    તથા $x \notin \emptyset$ની સ્પષ્ટ હકીકતો વપરાય છે.
  \end{quote}
  $\{x\}$નું એકમાત્ર ઘટક $x$ છે અને $x \in B$ છે;
  તેથી~$A$નું દરેક ઘટક~$B$નું ઘટક પણ છે.
  $\subseteq$ની વ્યાખ્યા મુજબ $A \subseteq B$.
\end{proof}

\subsection{અસ્તિત્વના દાવાનો ઉપયોગ}

ધારો કે અસ્તિત્વનો કોઈ દાવો સાચો છે તે તમે જાણો છો (કાં તો
તમે તેને સાબિત કર્યો છે અથવા તે વાપરી શકાય એવી પૂર્વધારણા છે),
જેમ કે ``કોઈ~$x$ માટે $x \in A$'' અથવા ``કોઈ $x \in A$
છે.'' સાબિતીમાં તેનો ઉપયોગ કરવો હોય તો પૂર્વધારણા મુજબ
અસ્તિત્વ ધરાવતી વસ્તુઓમાંથી એકને નામ આપ્યું છે એમ માની શકો.
$A$માં ઓછામાં ઓછી એક વસ્તુ છે, તેથી તે નામ જેનો ઉલ્લેખ
કરી શકે એવી વસ્તુઓ છે. તમે કદાચ તેમાંથી કોઈ ચોક્કસ વસ્તુ
પસંદ કરીને વધુ વર્ણવી ન શકો (તે $\in A$ છે એ સિવાય).
પરંતુ સાબિતી માટે તમે એક પસંદ કરી છે એમ માનીને તેને નામ
આપી શકો. અગાઉ વપરાયેલું અથવા પૂર્વધારણાઓમાં આવેલું નામ
ફરી ન વાપરવું અગત્યનું છે; નહીં તો ભૂલ થઈ શકે. સાબિતીમાં
``કોઈ $x$ માટે $x \in A$'' પરથી ``ધારો કે $a
\in A$'' તરફ જઈને આ દર્શાવો. હવે તમે~$a$ વિશે વિચાર કરી,
બીજી પૂર્વધારણાઓ વાપરીને વગેરે અંતે નિષ્કર્ષ $p$ સુધી પહોંચો.
જો $p$માં હવે~$a$નો ઉલ્લેખ ન હોય, તો $p$ ખાસ પસંદ કરેલા
$a \in A$ સાક્ષી પર આધાર રાખતું નથી; તે ``કોઈ $x$ માટે
$x \in A$'' એવી પૂર્વધારણાથી નીકળે છે.\sourcecorrection{OLMD-156}{મૂળમાં
નિષ્કર્ષને $a \in A$ની પૂર્વધારણાથી જ સ્વતંત્ર કહેવામાં આવ્યો છે.
અહીં ફક્ત ખાસ પસંદ કરેલા સાક્ષીથી સ્વતંત્રતા દર્શાવી છે;
અસ્તિત્વની પૂર્વધારણા હજી વપરાય છે.}

\begin{prop}
જો $A \neq \emptyset$, તો $A \cup B \neq \emptyset$.
\end{prop}

\begin{proof}
  $A \neq \emptyset$ ધારો. તેથી કોઈ $x$ માટે $x \in A$.
  \begin{quote}
    અહીં પહેલાં આપણે પ્રતિજ્ઞાપનની પૂર્વધારણા ફરી કહી છે. આ
    પૂર્વધારણા, એટલે $A \neq \emptyset$, અસ્તિત્વનો દાવો
    છુપાવે છે; કેટલીક વ્યાખ્યાઓ ઉકેલીએ ત્યારે તે મળે છે. $=$ની
    વ્યાખ્યા કહે છે કે $A = \emptyset$ ત્યારે અને ત્યારે જ જ્યારે
    દરેક $x \in A$ માટે $\in \emptyset$ પણ છે અને દરેક $x \in
    \emptyset$ માટે $\in A$ પણ છે. બંને બાજુનો નિષેધ કરતાં,
    $A \neq
    \emptyset$ ત્યારે અને ત્યારે જ જ્યારે કોઈ $x \in A$ માટે
    $\notin \emptyset$ છે અથવા કોઈ $x \in \emptyset$ માટે
    $\notin A$ છે. કોઈ વસ્તુ~$\in \emptyset$ નથી; તેથી બીજો
    વિકલ્પ ક્યારેય સાચો નથી, અને ``$x \in A$ અને $x
    \notin \emptyset$'' ફક્ત $x \in A$ બને છે. તેથી $A \neq
    \emptyset$ ત્યારે અને ત્યારે જ જ્યારે કોઈ $x$ માટે
    $x \in A$.\sourcecorrection{OLMD-157}{મૂળના આ વાક્યમાં
    ખોટી રીતે x અસમાન શૂન્યગણ લખાયું છે; સંદર્ભ અને અગાઉના
    પગલાં મુજબ A અસમાન શૂન્યગણ હોવું જોઈએ.} આ અસ્તિત્વનો
    દાવો છે. હવે~$A$ના ઘટકોમાંથી એકને નામ આપીને આ
    દાવો વાપરીએ:
  \end{quote}
  $a \in A$ લો.
  \begin{quote}
    હવે $\in A$ એવી એક વસ્તુને નામ આપ્યું છે. આગળ~$a$
    વિશે વિચારીશું, પરંતુ ફક્ત $a \in A$ જ ધારીશું, બીજું
    કંઈ નહીં:
  \end{quote}
  $a \in A$ હોવાથી $\cup$ની વ્યાખ્યા મુજબ $a \in A \cup B$.
  તેથી કોઈ $x$ માટે $x \in A \cup B$, એટલે $A \cup B \neq \emptyset$.
  \begin{quote}
    છેલ્લા પગલામાં ``$a \in A \cup B$'' પરથી ``કોઈ
    $x$ માટે $x \in A \cup B$'' મેળવ્યું. તેમાં હવે $a$નો
    ઉલ્લેખ નથી; એટલે ``કોઈ $x$ માટે $x \in A \cup B$''
    ફક્ત ``કોઈ $x$ માટે $x \in A$'' પરથી નીકળે છે. તેનો
    અર્થ $A \cup B \neq
    \emptyset$ છે.
  \end{quote}
\end{proof}

આપણે કર્યું તેમ ``$x$'' જેવા બંધિત ચલોને $a$ જેવાં કામચલાઉ
નામોથી જુદાં રાખવાનો અભ્યાસ સારો હોઈ શકે. પરંતુ વ્યવહારમાં
ઘણી વાર આમ કરતા નથી અને $x$ જ વાપરીએ છીએ, જેમ કે:
\begin{quote}
ધારો કે $A \neq \emptyset$, એટલે કોઈ $x \in A$ છે. $\cup$ની
વ્યાખ્યા મુજબ $x \in A \cup B$. તેથી $A \cup B \neq \emptyset$.
\end{quote}
પરંતુ આમ કરતાં જુદા અસ્તિત્વના દાવાઓ માટે જુદા $x$ અને $y$
વાપરો છો તેની ખાસ કાળજી રાખવી પડે. દાખલા તરીકે, નીચેનું
``જો $A \neq
\emptyset$ અને $B \neq \emptyset$ તો $A \cap B \neq \emptyset$''
સાબિત કરવાનું \emph{સાચું નથી} (કારણ કે એ વિધાન જ સાચું નથી):
\begin{quote}
ધારો કે $A \neq \emptyset$ અને $B \neq \emptyset$. તેથી કોઈ $x$ માટે $x
\in A$ અને કોઈ $x$ માટે $x \in B$. $x \in A$ અને $x \in
B$ હોવાથી $\cap$ની વ્યાખ્યા મુજબ $x \in A \cap B$. તેથી $A \cap B \neq
\emptyset$.
\end{quote}
ખોટું પગલું ક્યાં આવે છે અને પરિણામ કેમ સાચું નથી તે તમે
શોધી શકો છો?
% END OLP-0606
\endgroup


\begingroup
\makeatletter\def\ol@sectioncs{section}\makeatother

% BEGIN OLP-0607
\olfileid{mth}{prf}{ex1}

\olsection{એક ઉદાહરણ}
\phantomsection\label{guunit:OLP-0607}


આપણું પહેલું ઉદાહરણ ગણોના યોગગણ અને છેદગણ અંગેની નીચેની
સરળ હકીકત છે. તે વ્યાખ્યાઓ ઉકેલવી, સંયોજન અને સાર્વત્રિક દાવા
સાબિત કરવા તથા કિસ્સાવાર સાબિતી કરવી બતાવશે.

\begin{prop}
કોઈ પણ ગણો $A$, $B$ અને $C$ માટે $A \cup (B \cap C) = (A \cup B)
\cap (A \cup C)$.
\end{prop}

ચાલો તે સાબિત કરીએ!{}

\begin{proof}
આપણે બતાવવું છે કે કોઈ પણ ગણો $A$, $B$ અને $C$ માટે $A \cup (B \cap
C) = (A \cup B) \cap (A \cup C)$.
\begin{quote}
પહેલાં પ્રતિજ્ઞાપનમાંના ``$=$''ની વ્યાખ્યા ઉકેલીએ. યાદ રાખો કે
ગણો એકસમાન છે તે સાબિત કરવા તેમના ઘટકો સરખા છે તે
બતાવવું પડે. એટલે $A \cup
(B \cap C)$નાં બધાં ઘટકો, $(A \cup B) \cap (A \cup C)$નાં
પણ ઘટકો છે અને ઊલટું પણ છે. ``ઊલટું પણ''નો અર્થ
એ છે કે $(A
\cup B) \cap (A \cup C)$નું દરેક ઘટક, $A \cup (B \cap
C)$નું ઘટક પણ હોવું જોઈએ. તેથી વ્યાખ્યા ઉકેલતાં
આપણે સંયોજન સાબિત કરવાનું છે તે સમજાય છે. તેને નોંધીએ:
\end{quote}
વ્યાખ્યા મુજબ, $A \cup (B \cap C) = (A \cup B) \cap (A \cup C)$
ત્યારે અને ત્યારે જ જ્યારે $A \cup (B \cap C)$નું દરેક
ઘટક, $(A
\cup B) \cap (A \cup C)$નું ઘટક પણ છે અને
$(A \cup B) \cap (A
\cup C)$નું દરેક ઘટક, $A \cup (B \cap C)$નું
ઘટક પણ છે.
\begin{quote}
આ સંયોજન હોવાથી તેના બંને ભાગો અલગથી સાબિત કરવા પડે.
પહેલાથી શરૂ કરીએ: $A \cup (B \cap C)$નું દરેક
ઘટક, $(A
\cup B) \cap (A \cup C)$નું ઘટક પણ છે તે સાબિત કરીએ.

આ સાર્વત્રિક દાવો છે; તેથી $A \cup (B \cap C)$નું મનસ્વી ઘટક
લઈએ અને બતાવીએ કે તે $(A \cup B) \cap (A \cup C)$નું
ઘટક પણ હોવું જોઈએ. આ મનસ્વી ઘટકને કોઈ
ચલથી ઓળખાવીએ, ધારો કે~$z$. હવે સાબિતી આગળ વધે છે:
\end{quote}
પહેલાં આપણે સાબિત કરીએ કે $A \cup (B \cap C)$નું દરેક
ઘટક, $(A \cup B) \cap (A \cup C)$નું ઘટક
પણ છે. $z \in A \cup (B
\cap C)$ લો. આપણે $z \in (A \cup B) \cap (A \cup C)$ બતાવવું છે.
\begin{quote}
હવે $\cup$ અને~$\cap$ની વ્યાખ્યાઓ ઉકેલવાનો સમય છે. દાખલા તરીકે,
$\cup$ની વ્યાખ્યા $A \cup B = \Setabs{z}{z \in A
  \text{ અથવા } z \in B}$ છે. આ વ્યાખ્યા ``$A \cup (B
\cap C)$''ને લાગુ પાડીએ ત્યારે વ્યાખ્યાના ``$B$''નું સ્થાન
``$B \cap C$'' લે છે; તેથી $A \cup (B \cap C) = \Setabs{z}{z \in A \text{ અથવા }
  z \in B \cap C}$. એટલે $z \in A \cup (B \cap C)$
એવી આપણી પૂર્વધારણાનો અર્થ $z \in \Setabs{z}{z \in A \text{ અથવા } z \in B \cap
  C}$ થાય છે. અને $z \in \Setabs{z}{\dots z\dots}$ ત્યારે અને
ત્યારે જ જ્યારે \dots $z$ \dots; આ કિસ્સામાં $z \in A$
અથવા $z \in B \cap C$.
\end{quote}
$\cup$ની વ્યાખ્યા મુજબ, કાં તો $z \in A$ અથવા $z \in B \cap C$.
\begin{quote}
આ વિયોજન હોવાથી કિસ્સાવાર સાબિતી ઉપયોગી થશે. બે કિસ્સા લઈએ
અને દરેકમાં આપણો લક્ષિત નિષ્કર્ષ (એટલે ``$z \in (A \cup B) \cap (A \cup
C)$'') મળે છે તે બતાવીએ.
\end{quote}
કિસ્સો 1: ધારો કે $z \in A$.
\begin{quote}
આ પૂર્વધારણાઓમાંથી વધુ બહુ કંઈ કાઢવાનું નથી. એટલે નિષ્કર્ષમાં
આપણી પાસે શું છે તે જોઈએ. આપણે $z \in (A \cup B) \cap (A \cup C)$
બતાવવું છે. $\cap$ની વ્યાખ્યા મુજબ, $z \in (A \cup B) \cap (A \cup C)$
બતાવવા તે $(A \cup B)$ અને $(A \cup C)$ બંનેમાં છે તે બતાવવું
પડે. પરંતુ $z
\in A \cup B$ ત્યારે અને ત્યારે જ જ્યારે $z \in A$ અથવા $z \in B$;
અને કિસ્સો~1ની પૂર્વધારણા મુજબ $z \in A$ છે. એ જ
રીતે---$B$ના સ્થાને $C$ મૂકતાં---$z \in A \cup C$ મળે છે.
અહીં વિચાર ઊલટી દિશામાં થયો છે; તેથી સાબિતીમાં જરૂરી દિશામાં
આ તર્ક લખીએ.
\end{quote}
$z \in A$ હોવાથી $z \in A$ અથવા $z \in B$; એટલે~$\cup$ની
વ્યાખ્યા મુજબ $z \in A \cup B$. એ જ રીતે $z \in A \cup C$.
પરંતુ~$\cap$ની વ્યાખ્યા મુજબ તેનો અર્થ $z \in (A \cup B) \cap (A \cup C)$ થાય છે.
\begin{quote}
આથી કિસ્સાવાર સાબિતીનો પહેલો કિસ્સો પૂરો થયો. હવે બીજા
કિસ્સામાં, જ્યાં $z \in B \cap C$ છે, ત્યાં પણ નિષ્કર્ષ મેળવવો છે.
\end{quote}
કિસ્સો 2: ધારો કે $z \in B \cap C$.
\begin{quote}
ફરી આપણે બે ગણોના છેદગણ સાથે કામ કરીએ છીએ. $\cap$ની
વ્યાખ્યા લાગુ પાડીએ:
\end{quote}
$z \in B \cap C$ હોવાથી~$\cap$ની વ્યાખ્યા મુજબ $z$, $B$
અને $C$ બંનેનું ઘટક હોવું જોઈએ.
\begin{quote}
હવે ફરી નિષ્કર્ષ જોઈએ. $z$ એ $(A \cup B)$ અને $(A \cup C)$
બંનેમાં છે તે બતાવવું છે. ફરી જવાબ તરત મળે છે.
\end{quote}
$z \in B$ હોવાથી $z \in (A \cup B)$. $z \in C$ હોવાથી
$z \in (A
\cup C)$ પણ છે. તેથી $z \in (A \cup B) \cap (A \cup C)$.
\begin{quote}
અહીં ફરી $\cup$ અને $\cap$ની વ્યાખ્યાઓ લાગુ પાડી છે.
પરંતુ આપણે તે પહેલેથી યાદ કરી છે અને $z$ બેમાંથી કોઈ એક
ગણમાં હોય તો તેમના યોગગણમાં પણ છે તે અગાઉ બતાવ્યું છે;
તેથી આ પગલું એટલું સ્પષ્ટ લખવાની જરૂર નથી.

કિસ્સાવાર સાબિતીનો બીજો કિસ્સો પણ પૂરો થયો; હવે પહેલો
નિષ્કર્ષ કહી શકીએ.
\end{quote}
આથી જો $z \in A \cup (B \cap C)$, તો $z \in (A \cup B) \cap (A \cup C)$.
\begin{quote}
હવે વિપરીત દિશા બતાવવાની છે: $(A \cup B) \cap (A \cup C)$નું
દરેક ઘટક, $A \cup (B \cap
C)$નું ઘટક છે. પહેલાંની જેમ, પહેલા ગણનું મનસ્વી
ઘટક લઈને તે બીજા ગણમાં પણ છે તે બતાવીએ. હવે શું
કરવાના છીએ તે કહીએ.
\end{quote}
હવે ધારો કે $z \in (A \cup B) \cap (A \cup C)$. આપણે
$z \in A \cup (B \cap C)$ બતાવવું છે.
\begin{quote}
હવે $z \in (A \cup B) \cap (A
\cup C)$ એવી પૂર્વધારણાથી કામ કરીએ છીએ. સાબિતીના પહેલા ભાગ
જેવું જ~$z$ અહીં વાપરીએ છીએ, તેથી ગૂંચવણ ન થવી જોઈએ. પહેલો
ભાગ પૂરો કરતાં ત્યાંની બધી પૂર્વધારણાઓ હવે લાગુ નથી પડતી;
તેથી $z$ વિશે નવી પૂર્વધારણાઓ લઈ શકીએ. ગૂંચવણ થાય તો આગળના
ભાગમાં $z$ના સ્થાને બીજો ચલ વાપરો.

$\cap$ની વ્યાખ્યા મુજબ $z$, $A \cup B$ અને $A \cup C$
બંનેમાં છે. $\cup$ની વ્યાખ્યા ઉકેલતાં વધુમાં મળે છે: કાં તો
$z \in A$ અથવા $z \in B$, અને સાથે કાં તો $z \in A$
અથવા $z \in
C$. આ ફરી કિસ્સાવાર સાબિતી જેવું લાગે છે---પણ વચ્ચેનું
``અને'' ગૂંચવે છે. તમને લાગે કે માત્ર ત્રણ શક્યતાઓ છે:
$z$ એ $A$, $B$ કે $C$માંથી કોઈ એકમાં છે. પરંતુ તે ભૂલ
હશે. સાવચેતીથી દરેક વિયોજન અલગથી જોઈએ.
\end{quote}
$\cap$ની વ્યાખ્યા મુજબ $z \in A \cup B$ અને $z \in A \cup C$.
$\cup$ની વ્યાખ્યા મુજબ $z \in A$ અથવા $z \in B$.
કિસ્સા અલગ પાડીએ.
\begin{quote}
હાલમાં પહેલું વિયોજન લઈએ છીએ, તેથી $A \cup C$ ઉકેલીને
મળતું બીજું વિયોજન હજી વાપર્યું નથી. હકીકતમાં તેની હજી જરૂર
પણ નથી. પહેલો કિસ્સો $z \in A$ છે, અને કોઈ ગણનું
ઘટક તે ગણના બીજા કોઈ પણ ગણ સાથેના યોગગણનું પણ
ઘટક છે. તેથી કિસ્સો~1 સરળ છે:
\end{quote}
કિસ્સો 1: ધારો કે $z \in A$. તેથી $z \in A \cup (B \cap
C)$.
\begin{quote}
હવે બીજો કિસ્સો, $z \in B$. અહીં બીજો $\cup$ ઉકેલીને
ફરી કિસ્સા અલગ પાડીશું:
\end{quote}
કિસ્સો 2: ધારો કે $z \in B$. $z \in A \cup C$ હોવાથી
કાં તો $z \in
A$ અથવા $z \in C$. કિસ્સા વધુ અલગ પાડીએ:

કિસ્સો 2a: $z \in A$. તો ફરી $z \in A \cup (B \cap C)$.
\begin{quote}
આ થોડું વિચિત્ર હતું. આ કિસ્સામાં~$z \in B$ એવી
પૂર્વધારણાની ખરેખર જરૂર પડી નહીં; તેમાં વાંધો નથી.
\end{quote}
કિસ્સો 2b: $z \in C$. ત્યારે $z \in B$ અને $z \in C$;
તેથી $z \in B \cap C$ અને પરિણામે $z \in A \cup (B \cap C)$.
\begin{quote}
આથી બંને કિસ્સાવાર સાબિતીઓ પૂરી થાય છે અને વિપરીત દિશા
પણ સાબિત થાય છે.
\end{quote}
તેથી જો $z \in (A \cup B) \cap (A \cup C)$ તો $z \in A \cup (B \cap C)$.
\end{proof}
% END OLP-0607
\endgroup


\begingroup
\makeatletter\def\ol@sectioncs{section}\makeatother

% BEGIN OLP-0608
\olfileid{mth}{prf}{ex2}

\olsection{બીજું ઉદાહરણ}
\phantomsection\label{guunit:OLP-0608}


\begin{prop}
જો $A \subseteq C$, તો $A \cup (C \setminus A) = C$.
\end{prop}

\begin{proof}
  ધારો કે $A \subseteq C$. આપણે $A \cup (C
  \setminus A) = C$ બતાવવું છે.
  \begin{quote}
    પહેલાં નોંધીએ કે આ શરતી વિધાન છે. તેમાં સાર્વત્રિક પરિમાણક
    સ્પષ્ટ લખાયો નથી: પ્રતિજ્ઞાપન બધા ગણો $A$ અને $C$ માટે
    છે. તેથી $A$ અને $C$ મનસ્વી ગણો દર્શાવતા ચલો છે. આવું
    વિધાન સાબિત કરવા પૂર્વપક્ષ ધારીને પરિણામ સાબિત કરીએ.

    હવે $A \subseteq C$ની પૂર્વધારણા વાપરીએ. $\subseteq$ની
    વ્યાખ્યા ઉકેલીએ: તેનો અર્થ એ છે કે~$A$નાં બધાં
    ઘટકો, $C$નાં પણ ઘટકો છે. આ લખી
    લઈએ---આગળની આખી સાબિતીમાં આ હકીકત વાપરીશું.
  \end{quote}
  $\subseteq$ની વ્યાખ્યા મુજબ, $A \subseteq C$ હોવાથી દરેક
  $z$ માટે જો $z \in A$ તો $z \in C$.
  \begin{quote}
    પૂર્વધારણામાં આવેલી બધી વ્યાખ્યાઓ ઉકેલી લીધી. હવે
    નિષ્કર્ષ તરફ જઈએ. આપણે $A \cup (C \setminus A) = C$
    બતાવવું છે. તેથી અગાઉના ઉદાહરણની જેમ સાબિતી ગોઠવીએ:
    $A \cup (C \setminus A)$નું દરેક ઘટક~$C$નું
    ઘટક પણ છે અને ઊલટું, $C$નું દરેક ઘટક,
    $A \cup (C
    \setminus A)$નું ઘટક પણ છે તે બતાવીએ. તેને
    ટૂંકમાં $A \cup (C \setminus A)
    \subseteq C$ અને $C \subseteq A \cup (C \setminus A)$
    કહીએ. (અહીં વ્યાખ્યા ઉકેલવાને બદલે તેનો વિપરીત ઉપયોગ
    કરીએ છીએ, પણ તેથી સાબિતી વાંચવી થોડી સરળ બને છે.)
    આ સંયોજન હોવાથી બંને ભાગો સાબિત કરવા પડે. પહેલો ભાગ,
    એટલે $A \cup (C \setminus A)$નું દરેક ઘટક~$C$નું
    ઘટક પણ છે તે બતાવવા, મનસ્વી~$z$ માટે
    $z \in A \cup (C \setminus A)$ ધારીને $z \in C$
    બતાવીએ. $\cup$ની વ્યાખ્યા મુજબ $z \in A \cup (C \setminus A)$
    પરથી $z \in A$ અથવા $z \in C \setminus A$ મળે છે.
    હવે તમને આ રીત સમજાતી હશે.
    \end{quote}
  $A \cup (C \setminus A) = C$ ત્યારે અને ત્યારે જ જ્યારે
  $A \cup (C \setminus A) \subseteq
  C$ અને $C \subseteq (A \cup (C \setminus A))$ છે.\sourcecorrection{OLMD-158}{મૂળમાં
  અહીં જમણી બાજુના યોગગણના સૂત્રનું છેલ્લું કૌંસ છૂટી ગયું છે;
  અહીં તે પૂર્યું છે.} પહેલાં $A \cup (C \setminus A) \subseteq C$
  સાબિત કરીએ. $z \in A \cup (C
  \setminus A)$ લો. તેથી કાં તો $z \in A$ અથવા $z \in (C \setminus A)$.
  \begin{quote}
    આપણે વિયોજન સુધી પહોંચ્યા છીએ; તેમાંથી $z \in C$
    સાબિત કરવું છે. કિસ્સાવાર સાબિતી વાપરીએ.
  \end{quote}
  કિસ્સો 1: $z \in A$. દરેક $z$ માટે $z \in A$ હોય તો
  $z \in C$ હોવાથી અહીં $z \in C$ મળે છે.
  \begin{quote}
    અહીં પ્રતિજ્ઞાપનની $A \subseteq C$ પૂર્વધારણાથી મળેલી
    અગાઉ નોંધેલી હકીકત વાપરી છે. પહેલો કિસ્સો પૂરો થયો;
    હવે બીજા કિસ્સા $z \in (C
    \setminus A)$ તરફ જઈએ. યાદ રાખો કે $C \setminus A$ એ
    બે ગણોનો \emph{તફાવત ગણ} છે, એટલે~$C$નાં એ બધાં
    ઘટકોનો ગણ જે~$A$નાં ઘટકો નથી. પરંતુ
    $A$માં ન હોય એવું $C$નું કોઈ પણ ઘટક,~$C$નું
    ઘટક તો છે જ.
  \end{quote}
  કિસ્સો 2: $z \in (C \setminus A)$. એટલે $z \in C$ અને
  $z
  \notin A$. તેથી ખાસ કરીને $z \in C$.
  \begin{quote}
    સારું, પહેલી દિશા સાબિત થઈ. હવે બીજી દિશામાં
    $C \subseteq A \cup (C \setminus
    A)$ સાબિત કરીએ. એટલે $z \in C$ ધારીને $z \in A \cup (C
    \setminus A)$ બતાવીએ.
  \end{quote}
  હવે $z \in C$ લો. આપણે $z \in A$ અથવા $z \in C
  \setminus A$ બતાવવું છે.
  \begin{quote}
    $A$નાં બધાં ઘટકો, $C$નાં પણ ઘટકો છે;
    અને $C
    \setminus A$ એ $C$નાં એવા ઘટકોનો ગણ છે જે
    $A$માં નથી. તેથી $z$ કાં તો $A$માં છે અથવા $C
    \setminus A$માં છે. પરિણામ પહેલેથી સાચું છે તે ખબર ન હોય
    તો આ પગલું અસ્પષ્ટ લાગશે. તેને પગલાંવાર સાબિત કરવું સારું.
    સાબિતી વિના કહી શકાય એવી સરળ હકીકત મદદ કરશે:
    $z \in A$ અથવા $z \notin A$. આને ``વર્જિત મધ્યનો નિયમ''
    કહે છે: કોઈ પણ વિધાન~$p$ માટે કાં તો $p$ સાચું છે અથવા
    તેનો નિષેધ સાચો છે. (અહીં $p$ એટલે $z \in A$ એ વિધાન.)
    આ વિયોજન હોવાથી ફરી કિસ્સાવાર સાબિતી કરી શકીએ.
  \end{quote}
  કાં તો $z \in A$ અથવા $z \notin A$. પહેલા કિસ્સામાં
  $z \in A \cup
  (C \setminus A)$. બીજા કિસ્સામાં $z \in C$ અને
  $z \notin A$, તેથી $z \in C \setminus A$. પરિણામે
  $z \in A \cup (C \setminus A)$.
  \begin{quote}
    સાબિતી પૂરી થઈ: આપણે $A \cup (C \setminus A) = C$ બતાવ્યું છે.
  \end{quote}
\end{proof}
% END OLP-0608
\endgroup


\begingroup
\makeatletter\def\ol@sectioncs{section}\makeatother

% BEGIN OLP-0609
\olfileid{mth}{prf}{con}

\olsection{વિરોધાભાસથી સાબિતી}
\phantomsection\label{guunit:OLP-0609}


સૌપ્રથમ, વિરોધાભાસથી સાબિતી નકારાત્મક દાવા સાબિત કરવા વપરાતો
અનુમાનનો નમૂનો છે. ધારો કે કોઈ દાવો~$p$ \emph{ખોટો} છે,
એટલે~$\lnot p$ છે, તે બતાવવું છે. સૌથી આશાસ્પદ રીત આ છે:
(a) $p$ સાચું છે એમ ધારો અને (b) બતાવો કે આ પૂર્વધારણાથી
ખોટી હોવાનું જાણીતું કંઈક નીકળે છે. ``ખોટી હોવાનું જાણીતું''
પરિણામ પોતે~$p$ સાથે અથવા તપાસાતા સમગ્ર દાવાની બીજી કોઈ
પૂર્વધારણા સાથે ટકરાઈ શકે. દાખલા તરીકે, ``જો $q$ તો $\lnot p$''
સાબિત કરવા $q$ સાચું ધારીને તેમાંથી~$\lnot
p$ સાબિત કરવું પડે. જો $\lnot p$ વિરોધાભાસથી સાબિત કરો,
તો~$q$ ઉપરાંત $p$ પણ ધારવું પડે. $p$ પરથી $\lnot q$
સાબિત થાય, તો~$p$ની પૂર્વધારણાથી તમારી બીજી પૂર્વધારણા~$q$
સાથે વિરોધાભાસ મળે છે, કારણ કે $q$ અને $\lnot q$ બંને
સાચાં હોઈ શકતાં નથી. અલબત્ત, વિરોધાભાસ મેળવવા સાબિતીમાં
બીજા અનુમાન-નમૂનાઓ અને વ્યાખ્યાઓ ઉકેલવાની રીત પણ વાપરવી
પડે. એક ઉદાહરણ જોઈએ.

\begin{prop}
  જો $A \subseteq B$ અને $B = \emptyset$, તો $A$માં કોઈ
  ઘટકો નથી.
\end{prop}

\begin{proof}
  $A \subseteq B$ અને $B = \emptyset$ ધારો. આપણે
  $A$માં કોઈ ઘટકો નથી તે બતાવવું છે.
  \begin{quote}
    આ શરતી દાવો હોવાથી પૂર્વપક્ષ ધારીને પરિણામ સાબિત
    કરવાનું છે. પરિણામ છે: $A$માં કોઈ ઘટકો નથી.
    વધુ સ્પષ્ટ કહીએ તો એવો~$x \in A$ છે એવું નથી.
  \end{quote}
  $A$માં કોઈ ઘટકો નથી ત્યારે અને ત્યારે જ જ્યારે
  એવો~$x$ છે કે $x \in A$ એવું નથી.
  \begin{quote}
    એટલે આપણે ખરેખર નકારાત્મક દાવો~$\lnot p$ સાબિત કરવો છે:
    એવો $x \in A$ છે એવું નથી. વિરોધાભાસથી સાબિતી કરવા
    અનુરૂપ હકારાત્મક દાવો~$p$, એટલે એવો $x \in A$ છે, એમ
    ધારીને તેમાંથી વિરોધાભાસ મેળવવો પડે. ``વિરોધાભાસ માટે
    ધારો કે'' અથવા ફક્ત ``ધારો કે એવું નથી'' લખીને આ રીત
    દર્શાવીએ છીએ અને પછી પૂર્વધારણા~$p$ કહીએ છીએ.
  \end{quote}
  ધારો કે એવું નથી: કોઈ $x \in A$ છે.
  \begin{quote}
    હવે આ નવી પૂર્વધારણાથી વિરોધાભાસ મેળવવો છે. બીજી બે
    પૂર્વધારણાઓ $A \subseteq B$ અને $B = \emptyset$ છે.
    પહેલી પરથી $x \in B$ મળે છે:
  \end{quote}
  $A \subseteq B$ હોવાથી $x \in B$.
  \begin{quote}
    પરંતુ $B = \emptyset$ હોવાથી $B$નું દરેક ઘટક
    (જેમ કે $x$)~$\emptyset$નું પણ ઘટક હોવું જોઈએ.
  \end{quote}
  $B = \emptyset$ હોવાથી $x \in \emptyset$. આ વિરોધાભાસ છે,
  કારણ કે વ્યાખ્યા મુજબ $\emptyset$માં કોઈ ઘટકો નથી.
  \begin{quote}
    આથી સાબિતી પૂરી થાય છે: લીધેલી પૂર્વધારણાઓમાંથી જરૂરી
    વિરોધાભાસ મળ્યો છે, તેથી બધી પૂર્વધારણાઓ એકસાથે સાચી
    હોઈ શકતી નથી. પહેલી બે પૂર્વધારણાઓ ($A \subseteq B$ અને
    $B = \emptyset$) પર શંકા નથી, તેથી છેલ્લે લીધેલી
    પૂર્વધારણા (એવો $x \in A$ છે) ખોટી હોવી જોઈએ. વિગત
    જોઈએ તો તેને સ્પષ્ટ લખી શકીએ.
  \end{quote}
  તેથી એવો $x \in A$ છે એવી આપણી પૂર્વધારણા ખોટી છે;
  એટલે વિરોધાભાસથી સાબિતી મુજબ $A$માં કોઈ ઘટકો નથી.
\end{proof}

શાસ્ત્રીય તર્કમાં દરેક હકારાત્મક દાવો નકારાત્મક દાવાને સમકક્ષ
છે: $p$ ત્યારે અને ત્યારે જ જ્યારે $\lnot\lnot p$.
\sourcecorrection{OLMD-159}{મૂળમાં આ દ્વિનિષેધ સમકક્ષતા કોઈ
વ્યાપશરત વિના કહેવાઈ છે. તે શાસ્ત્રીય તર્કમાં માન્ય છે; સ્ફુરણાવાદી
તર્કમાં સામાન્ય રીતે નહીં.} તેથી વિરોધાભાસથી સાબિતી હકારાત્મક
દાવાને ``પરોક્ષ રીતે'' સ્થાપિત કરવા પણ વાપરી શકાય છે:~$p$
સાબિત કરવા તેને નકારાત્મક દાવો $\lnot\lnot p$ તરીકે વાંચો.
$\lnot p$ પરથી વિરોધાભાસ સાબિત થાય, તો વિરોધાભાસથી સાબિતી
મુજબ $\lnot\lnot p$ મળે અને તેથી~$p$ મળે.

છેલ્લા ઉદાહરણમાં આપણું લક્ષ્ય $A$માં કોઈ ઘટકો નથી એવો
નકારાત્મક દાવો હતો. તેથી વિરોધાભાસથી સાબિતી માટે લીધેલી
પૂર્વધારણા (એવો $x \in A$ છે) હકારાત્મક હતી. તેમાંથી વિચારવા
માટે કશુંક મળ્યું, એટલે કામચલાઉ $x \in A$; તેના વિશે આગળ
વિચારતાં $x \in \emptyset$ સુધી પહોંચ્યા.

હકારાત્મક દાવો પરોક્ષ રીતે સાબિત કરતાં વિરોધાભાસ માટેની
પૂર્વધારણા નકારાત્મક હોય છે. પરંતુ ઘણી વાર હકારાત્મક દાવાને
નકારાત્મક સ્વરૂપમાં અને નકારાત્મક દાવાને હકારાત્મક સ્વરૂપમાં
સહેલાઈથી કહી શકાય. જો છેલ્લા ઉદાહરણમાં ``$A$માં કોઈ
ઘટકો નથી'' એ નકારાત્મક પરિણામને બદલે ``$A = \emptyset$''
સાબિત કર્યું હોત તો સાબિતી લગભગ એવી જ હોત. ($=$ની વ્યાખ્યા
મુજબ ``$A = \emptyset$'' સાર્વત્રિક દાવો છે: તેને ઉકેલતાં
``~$A$નું દરેક ઘટક,~$\emptyset$નું ઘટક છે
અને ઊલટું પણ'' મળે છે.) પરંતુ તે ``એવો $x \in A$ છે
એવું નથી'' એ નકારાત્મક દાવાને સમકક્ષ છે તે સહેલાઈથી જોઈ શકાય.

આથી ક્યારેક~$p$ને સીધું સાબિત કરવા કરતાં $\lnot p$ને
પૂર્વધારણા તરીકે વાપરવું સરળ હોય છે. સીધી સાબિતી એટલી જ સરળ
અથવા વધુ સરળ હોય (જેમ કે આગળનાં ઉદાહરણોમાં), તોપણ કેટલાક
લોકો પરોક્ષ રીતે આગળ વધવાનું પસંદ કરે છે. દ્વિનિષેધ ગૂંચવે તો
વિરોધાભાસથી કોઈ દાવાની સાબિતીને તેના \emph{વિપરીત} દાવામાંથી
વિરોધાભાસ મેળવવાની સાબિતી માનો. એટલે $\lnot
p$ની વિરોધાભાસથી સાબિતી એટલે~$p$ ધારીને વિરોધાભાસ મેળવવો;
અને~$p$ની વિરોધાભાસથી સાબિતી એટલે~$\lnot p$ ધારીને
વિરોધાભાસ મેળવવો.

\begin{prop}
$A \subseteq A \cup B$.
\end{prop}

\begin{proof}
  આપણે $A \subseteq A \cup B$ બતાવવું છે.
  \begin{quote}
    ઉપરથી આ હકારાત્મક દાવો છે: દરેક $x \in A$, $A \cup B$માં
    પણ છે. તેનો નિષેધ છે: કોઈ $x \in A$ એવો છે જે
    $\notin A \cup B$ છે. આ નકારેલા દાવાને ધારીને તેમાંથી
    વિરોધાભાસ મેળવી દાવો પરોક્ષ રીતે સાબિત કરી શકીએ.
    \end{quote}
  ધારો કે એવું નથી, એટલે $A \nsubseteq A \cup B$.
  \begin{quote}
    $A \subseteq A \cup B$ની વ્યાખ્યા આપણી પાસે છે: દરેક
    $x \in A$, $\in A \cup B$ પણ છે. $A \nsubseteq A \cup B$
    શું કહે છે તે સમજવા ઉકેલેલી વ્યાખ્યા પર થોડું મૂળભૂત તાર્કિક
    કામ કરવું પડે: દરેક $x \in A$~$\in
    A \cup B$ પણ છે એવું ખોટું હોય ત્યારે અને ત્યારે જ જ્યારે
    \emph{કોઈ}~$x \in A$ એવો છે જે $\notin A \cup B$ છે.
    \sourcecorrection{OLMD-160}{મૂળના આ વાક્યમાં અહીં અચાનક
    Cમાં સભ્ય નથી લખાયું છે; દાવાના નિષેધ મુજબ A યોગ Bમાં
    સભ્ય નથી હોવું જોઈએ.} (અહીં ખાસ સાવધ રહો: ``બધા $A$s
    એ $B$s છે'' જેવા દાવાનો નિષેધ ખોટી રીતે સમજવાથી
    વિદ્યાર્થીઓની ઘણી વિરોધાભાસી સાબિતીઓ ખોટી પડે છે.)
    બીજા શબ્દોમાં, $A \nsubseteq A \cup B$ ત્યારે અને ત્યારે જ
    જ્યારે એવો $x$ છે કે $x \in A$ અને $x \notin A \cup B$.
    પછીનું સરળ છે.
  \end{quote}
  તેથી એવો $x \in A$ છે કે $x \notin A \cup B$. $\cup$ની
  વ્યાખ્યા મુજબ $x \in A \cup B$ ત્યારે અને ત્યારે જ જ્યારે
  $x \in A$ અથવા $x \in
  B$. $x \in A$ હોવાથી $x \in A \cup B$ મળે છે. પરંતુ આ
  $x \notin A \cup B$ની પૂર્વધારણાથી વિરોધાભાસી છે.
\end{proof}

\begin{prob}
$A \cap B \subseteq A$ \emph{પરોક્ષ રીતે} સાબિત કરો.
\end{prob}

\begin{prop}
જો $A \subseteq B$ અને $B \subseteq C$, તો $A \subseteq C$.
\end{prop}

\begin{proof}
  $A \subseteq B$ અને $B \subseteq C$ ધારો. આપણે $A
  \subseteq C$ બતાવવું છે.
  \begin{quote}
    પરોક્ષ રીતે આગળ વધીએ: જે સાબિત કરવું છે તેનો નિષેધ ધારો.
  \end{quote}
  ધારો કે એવું નથી, એટલે $A \nsubseteq C$.
  \begin{quote}
    અગાઉની જેમ, $A \nsubseteq C$ ત્યારે અને ત્યારે જ જ્યારે
    દરેક $x \in A$, $\in C$ પણ છે એવું નથી; એટલે કોઈ
    $x \in A$ એવો છે જે $\notin C$ છે. ચિંતા ન કરો,
    અભ્યાસથી આવા નિષેધ ઉકેલવા વધારે વિચારવું નહીં પડે.
  \end{quote}
  બીજા શબ્દોમાં, એવો~$x$ છે કે $x \in A$ અને $x \notin C$.
  \begin{quote}
    હવે અહીંથી વિરોધાભાસ મેળવીએ. એ માટે બીજી બે પૂર્વધારણાઓ
    પણ વાપરવી પડશે.
  \end{quote}
  $A \subseteq B$ હોવાથી $x \in B$. $B \subseteq C$ હોવાથી
  $x \in
  C$. પરંતુ આ $x \notin C$ સાથે વિરોધાભાસ છે.
\end{proof}

\begin{prop}
જો $A \cup B = A \cap B$ તો $A = B$.
\end{prop}

\begin{proof}
  $A \cup B = A \cap B$ ધારો. આપણે $A = B$ બતાવવું છે.
  \begin{quote}
    હવે શરૂઆતની રીત જાણીતી છે:
  \end{quote}
  વિરોધાભાસ માટે $A \neq B$ ધારો.
  \begin{quote}
    વિરોધાભાસી સાબિતીની આપણી પૂર્વધારણા $A \neq
    B$ છે. $A = B$ ત્યારે અને ત્યારે જ જ્યારે $A \subseteq B$
    અને $B \subseteq A$; તેથી $A \neq B$ ત્યારે અને ત્યારે જ
    જ્યારે $A \nsubseteq B$ \emph{અથવા} $B \nsubseteq
    A$. (નિષેધનો વિચાર કરતી વખતે સાવધ રહેવું કેટલું અગત્યનું
    છે તે જુઓ!{}) આ વિયોજનમાંથી વિરોધાભાસ મેળવવા કિસ્સાવાર
    સાબિતી કરીએ: દરેક કિસ્સામાં વિરોધાભાસ મળે છે તે બતાવીએ.
  \end{quote}
  $A \neq B$ ત્યારે અને ત્યારે જ જ્યારે $A \nsubseteq B$ અથવા
  $B \nsubseteq A$. કિસ્સા અલગ પાડીએ.
  \begin{quote}
    પહેલા કિસ્સામાં $A \nsubseteq B$ ધારો; એટલે કોઈ $x$
    માટે $x \in A$ પણ $\notin B$. $A \cap B$ એ $A$ અને
    $B$ બંનેનાં સામાન્ય ઘટકોનો ગણ છે; તેથી કોઈ
    વસ્તુ એક ગણમાં ન હોય તો તે છેદગણમાં પણ નથી. $A \cup B$
    એ $A$ અને $B$નાં બધા સભ્યોનો ગણ છે, એટલે કોઈ એકમાં
    હોય તે યોગગણમાં પણ છે. તેથી $x \in A \cup B$ પણ
    $x \notin A \cap
    B$, અને આથી $A \cap B \neq A \cup B$.
  \end{quote}

  કિસ્સો 1: $A \nsubseteq B$. તો કોઈ $x$ માટે $x \in A$
  પણ $x \notin
  B$. $x \notin B$ હોવાથી $x \notin A \cap B$. $x \in A$
  હોવાથી $x \in A \cup B$. તેથી $A \cap B \neq A \cup B$,
  જે $A \cap B = A \cup B$ની પૂર્વધારણાથી વિરોધાભાસી છે.

  કિસ્સો 2: $B \nsubseteq A$. તો કોઈ $y$ માટે $y \in B$
  પણ $y \notin
  A$. અગાઉની જેમ $y \in A \cup B$ પરંતુ $y \notin A \cap B$;
  તેથી $A \cap B \neq A \cup B$, જે ફરી $A \cap B = A \cup
  B$થી વિરોધાભાસી છે.
\end{proof}
% END OLP-0609
\endgroup


\begingroup
\makeatletter\def\ol@sectioncs{section}\makeatother

% BEGIN OLP-0610
\olfileid{mth}{prf}{rea}
\olsection{સાબિતીઓ વાંચવી}
\phantomsection\label{guunit:OLP-0610}


પાઠ્યપુસ્તકો અને લેખોમાં મળતી સાબિતીઓમાં અત્યાર સુધીના
ઉદાહરણો જેવી બધી વિગતો ભાગ્યે જ લખાય છે. લેખકો ઘણી વાર
કિસ્સા અલગ પાડે છે, પરોક્ષ સાબિતી કરે છે, કે વ્યાખ્યા વાપરે છે
તે તરફ ધ્યાન દોરતા નથી. તેથી પાઠ્યપુસ્તકની સાબિતી વાંચતાં
તે સમજવા આવી વિગતો તમારે જાતે પૂરી કરવી પડે છે. આમ કરવાથી
સાબિતીમાં કરવાનાં વિવિધ પગલાંનો અભ્યાસ પણ થાય છે. એક
ઉદાહરણ જોઈએ.

\begin{prop}[શોષણ]
બધા ગણો $A$, $B$ માટે,
\[
A \cap (A \cup B) = A
\]
\end{prop}

\begin{proof}
જો $z \in A \cap (A \cup B)$, તો $z \in A$; તેથી $A \cap (A \cup B)
\subseteq A$. હવે $z \in A$ ધારો. ત્યારે $z \in A \cup B$ પણ છે;
તેથી $z \in A \cap (A \cup B)$ પણ છે.
\end{proof}

ઉપર શોષણના નિયમની સાબિતી ખૂબ સંક્ષિપ્ત છે. કઈ વ્યાખ્યાઓ
વપરાઈ છે તેનો ઉલ્લેખ નથી, સાબિત કરતા પહેલાં ``આ બતાવવાનું
છે'' એવું પણ લખાયું નથી. તેને ઉકેલીએ. સાબિત પ્રતિજ્ઞાપન
મનસ્વી ગણો $A$ અને $B$ અંગેનો સાર્વત્રિક દાવો છે; સાબિતીમાં
$A$ કે $B$ આવે ત્યારે તે મનસ્વી ગણો દર્શાવતા ચલો છે.
સાબિતી જે સામાન્ય દાવા સ્થાપિત કરે છે તે ગણોની સમાનતા
સાબિત કરવા જરૂરી છે: સમાનતાની ડાબી બાજુનું દરેક ઘટક
જમણી બાજુનું ઘટક છે અને ઊલટું પણ.

\begin{quote}
``જો $z \in A \cap (A \cup B)$, તો $z \in A$; તેથી $A \cap (A \cup B)
  \subseteq A$.''
\end{quote}

આ સમાનતાની સાબિતીનો પહેલો ભાગ છે: ડાબી બાજુનું મનસ્વી~$z$
ઘટક હોય, તો તે જમણી બાજુનું પણ ઘટક છે,
એટલે $A \cap (A \cup B) \subseteq A$. $z \in A \cap (A \cup B)$
ધારો. બે ગણોના છેદગણમાં $z$ ઘટક હોય ત્યારે અને ત્યારે જ
જ્યારે તે બંને ગણોનું ઘટક હોય છે; તેથી $z \in A$
અને $z \in A \cup B$ મળે. ખાસ કરીને $z \in A$, જે બતાવવું
હતું. પહેલા ભાગ માટે આટલું પૂરતું હોવાથી બાકીની સાબિતી
બીજા ભાગની છે: $A \subseteq A \cap
(A \cup B)$ સાબિત કરવાની.

\begin{quote}
``હવે $z \in A$ ધારો. ત્યારે $z \in A \cup B$ પણ છે;
તેથી $z \in A \cap (A \cup B)$ પણ છે.''
\end{quote}

દરેક~$z$ માટે જો $z \in A$ તો $z \in A \cap (A \cup B)$
બતાવીએ છીએ, તેથી $z \in A$ ધારીને શરૂ કરીએ. $z
\in A \cap (A \cup B)$ બતાવવા $\cap$ની વ્યાખ્યા મુજબ
(i) $z \in A$ અને (ii) $z \in A \cup B$ બંને બતાવવું પડે.
અહીં (i) આપણી પૂર્વધારણા જ છે; વધુ કંઈ સાબિત કરવાનું નથી.
એટલા માટે સાબિતીમાં તેનો ફરી ઉલ્લેખ નથી. (ii) માટે યાદ કરો:
$z$ ગણોના યોગગણનું ઘટક ત્યારે અને ત્યારે જ જ્યારે
તેમાંના ઓછામાં ઓછા એક ગણનું ઘટક હોય. $z \in A$
છે અને $A \cup B$ એ $A$ અને $B$નો યોગગણ છે, એટલે અહીં
આ શરત પૂરી થાય છે. તેથી $z \in A \cup B$. (i) $z \in A$
અને (ii) $z \in A \cup B$ બંને બતાવ્યાં; એટલે $\cap$ની
વ્યાખ્યા મુજબ $z \in A \cap (A \cup B)$. સાબિતી આ
વ્યાખ્યાઓનો ઉલ્લેખ કરતી નથી, કારણ કે વાચકને તે આવડી ગઈ છે
એમ માનવામાં આવે છે. ન આવડી હોય તો પાછા જઈને યાદ કરો.
પછી $z
\in A$ પરથી $z \in A \cup B$ કેમ મળે છે અને $z \in A$
તથા $z \in A \cup B$ પરથી $z \in A \cap (A \cup B)$
કેમ મળે છે તે પણ ઓળખવું પડે.

ઉપરની સાબિતીનું બધું સ્પષ્ટ કરતું બીજું સ્વરૂપ આ છે:
\begin{proof}{}
[$=$ની ગણો માટેની વ્યાખ્યા મુજબ $A \cap (A \cup B) = A$
બતાવવા (a) $A \cap (A \cup B) \subseteq A$ અને
(b) $A \subseteq A \cap (A \cup B)$ બતાવવું પડે.\sourcecorrection{OLMD-161}{મૂળના
(b)માં ભૂલથી (a) જેવી જ ઉપગણ-દિશા ફરી લખાઈ છે;
અહીં વિપરીત દિશા પૂરી છે.} (a): $\subseteq$ની વ્યાખ્યા
મુજબ જો $z \in A \cap (A \cup B)$ તો $z \in A$ બતાવવું પડે.]
જો $z \in A
\cap (A \cup B)$, તો $z \in A$ [$\cap$ની વ્યાખ્યા મુજબ
$z
  \in A \cap (A \cup B)$ ત્યારે અને ત્યારે જ જ્યારે
$z \in A$ અને $z \in A \cup B$], તેથી $A
\cap (A \cup B) \subseteq A$. [(b): $\subseteq$ની વ્યાખ્યા
મુજબ જો $z \in A$, તો $z \in A \cap (A \cup B)$ બતાવવું પડે.]
હવે [(1)] $z \in A$ ધારો. ત્યારે [(2)] $z \in A \cup B$ પણ છે
[(1) પરથી $z \in A$ અથવા $z \in B$, અને $\cup$ની વ્યાખ્યા
મુજબ તેનો અર્થ $z
  \in A \cup B$ થાય છે]; તેથી $z \in A \cap (A \cup B)$ પણ
છે [$\cap$ની વ્યાખ્યા મુજબ $z \in A$, એટલે (1), અને
$z
  \in A \cup B$, એટલે (2), બંને જોઈએ].\sourcecorrection{OLMD-162}{મૂળમાં
છેલ્લા યોગગણ-સભ્યતા સૂત્રમાં વધારાનું બંધ કૌંસ છે; અહીં
તે દૂર કર્યું છે.}
\end{proof}

\begin{prob}
$A \cup (A \cap B) = A$ની નીચેની સાબિતીને વિસ્તારો:
વપરાયેલા બધા અનુમાન-નમૂનાઓ જણાવો, દરેક પગલું પૂર્વધારણા
અથવા અગાઉ સ્થાપિત દાવામાંથી કેમ નીકળે છે તે કહો, અને
ક્યાં કઈ વ્યાખ્યાનો ઉપયોગ થાય છે તે બતાવો.
\begin{proof}
  જો $z \in A \cup (A \cap B)$, તો $z \in A$ અથવા
  $z \in A \cap B$. જો $z \in A \cap B$, તો $z \in A$.
  કોઈ પણ $z \in A$ એ $\in A \cup (A
  \cap B)$ પણ છે.
\end{proof}
\end{prob}
% END OLP-0610
\endgroup


\begingroup
\makeatletter\def\ol@sectioncs{section}\makeatother

% BEGIN OLP-0611
\olfileid{mth}{prf}{cnt}

\olsection{મારાથી આ નહીં થાય!{}}
\phantomsection\label{guunit:OLP-0611}


આપણને સૌને ક્યારેક હાર માની લેવાનું મન થાય છે. પણ તમે તે
\emph{કરી શકો છો}. તમારા અધ્યાપક અને શિક્ષણ-સહાયક તેમજ
સહાધ્યાયીઓ મદદ કરી શકે છે. તેમની મદદ માગો!{} મુશ્કેલી
ટાળવા અને હાર માનવાનું મન થાય ત્યારે શું કરવું તે માટે
આ કેટલીક સૂચનાઓ છે.

પ્રશ્નો સફળતાપૂર્વક ઉકેલી શકો તે માટે આ કરો:
\begin{enumerate}
\item \emph{શક્ય હોય એટલું વહેલું શરૂ કરો.} સત્ર દરમિયાન
  આપણે સૌ વ્યસ્ત થઈ જઈએ છીએ અને ઘણાને કામ ટાળવાની ટેવ
  સતાવે છે. તેથી ઘરકામ વહેલું શરૂ કરવું શ્રેષ્ઠ ઉપાયોમાંનો
  એક છે. એમ કરશો તો અટકો ત્યારે ઉકેલ શોધવા માટે સમય
  રહેશે (રડી પડવું ઉકેલ નથી).
\item \emph{સહાધ્યાયીઓ સાથે વાત કરો}. તમે એકલા નથી.
  વર્ગમાં બીજા લોકોને પણ મુશ્કેલી પડતી હશે---ભલે તેમની
  મુશ્કેલી જુદી હોય. તેમની સાથે વાત કરવાથી પ્રશ્નનો બીજો
  દૃષ્ટિકોણ મળે અને ઉકેલ સૂઝી શકે. અલબત્ત, તેમનો ઉકેલ
  સીધો ન ઉતારો: સંકેત માગો, અથવા તમે ક્યાં અટક્યા છો તે
  સમજાવીને આગળનું પગલું પૂછો. પછી તમને સમજ પડે ત્યારે
  તેમને પણ મદદ કરો. બીજાને સમજાવવાથી તમારી પોતાની
  સમજ પણ વધુ સ્પષ્ટ થાય છે.
\item \emph{મદદ માગો.} તમારા માટે ઘણા સ્રોતો છે---તમારા
  અધ્યાપક અને શિક્ષણ-સહાયક તમારી મદદ માટે છે અને તમે
  સફળ થાઓ એવું \emph{ઇચ્છે} છે. પ્રશ્ન ઉકેલવામાં અને
  કયા તબક્કે અટકો છો તે ઓળખવામાં તેઓ મદદ કરી શકશે.
\item \emph{થોડો વિરામ લો.} અટક્યા હો તો તેનું એક કારણ
  \emph{કદાચ} એ હોય કે તમે પ્રશ્નને બહુ લાંબો સમયથી જોઈ રહ્યા છો.
  નાનો વિરામ લો, ચા પીવો અથવા થોડો સમય બીજો પ્રશ્ન
  ઉકેલો; પછી તાજા મનથી પાછા આવો. એક રાત ઊંઘ્યા પછી
  ફરી વિચારો.
\end{enumerate}

ધ્યાન આપ્યું કે આ રીતો માટે સાબિતી પર ઘણું પહેલાંથી કામ
શરૂ કરેલું હોવું જોઈએ? સોંપવાની તારીખના આગલા દિવસે
વહેલી સવારે બે વાગ્યે શરૂ કર્યું હોય તો કદાચ તે બહુ
ઉપયોગી ન થાય.

આ નિરાશાજનક લાગશે, પણ સાબિતીનો પ્રશ્ન ઉકેલવો એવો પડકાર
છે જે અંતે સંતોષ આપે છે. કેટલાક લોકો આને વ્યવસાય તરીકે
કરે છે---એટલે તેમાં આનંદ લેવા જેવું કંઈક હશે જ. લગભગ
બીજી દરેક બાબતની જેમ પ્રશ્નો ઉકેલવા અને સાબિતીઓ કરવી
અભ્યાસ માગે છે. કેટલાક સહાધ્યાયીઓને આ સરળ લાગતું હશે;
તેમણે કદાચ પહેલેથી ઘણો અભ્યાસ કર્યો હશે. બહુ સહેલાઈથી
હાર ન માનો.

કોઈ પ્રશ્ન માટે સમય (અથવા ધીરજ) ખૂટી જાય તો તેમાં વાંધો
નથી. તેનો અર્થ તમે બુદ્ધિહીન છો કે તમને કદી નહીં સમજાય
એવો નથી. અધ્યાપક અથવા બીજા વિદ્યાર્થી પાસેથી ઉકેલની
રીત જાણો અને તમે ક્યાં ભૂલ્યા અથવા અટક્યા તે ઓળખો,
જેથી આગળ આવો પ્રશ્ન આવે ત્યારે એ ભૂલ ટાળી શકો. પછી
ઉકેલ જોયા વિના ફરી જાતે કરવાનો પ્રયાસ કરો. અને આગલી
વખતે વહેલું શરૂ કરો (અને વહેલું મદદ માગો).
% END OLP-0611
\endgroup


\begingroup
\makeatletter\def\ol@sectioncs{section}\makeatother

% BEGIN OLP-0612
\olfileid{mth}{prf}{res}

\section{અન્ય સ્રોતો}
\phantomsection\label{guunit:OLP-0612}


ગણિતમાં સાબિતીઓ કેવી રીતે કરવી તે અંગે ઘણાં ઉપયોગી પુસ્તકો
છે. ખાસ કરીને \emph{How to Read and do Proofs: An Introduction to
Mathematical Thought Processes} \citep{Solow2013} અને \emph{How to
Prove It: A Structured Approach} \citep{Velleman2019} જુઓ.
\href{http://www.people.vcu.edu/~rhammack/BookOfProof/BookOfProof.pdf}{\emph{Book
of Proof}} \citep{Hammack2013} અને
\href{https://scholarworks.gvsu.edu/books/7/}{\emph{Mathematical
Reasoning}} \citep{Sandstrum2019} એ સાબિતી અંગેનાં પુસ્તકો
ઓનલાઇન નિઃશુલ્ક ઉપલબ્ધ છે. તત્ત્વચિંતકોને ગાણિતિક તર્કવિચારના
પ્રારંભિક પરિચય માટે
\emph{More
Precisely: The Math you need to do Philosophy} \citep{Steinhart2018}
ઉપયોગી લાગી શકે.

ઇન્ટરનેટ પર સાબિતી અંગે ટૂંકી માર્ગદર્શિકાઓ પણ છે;
દાખલા તરીકે,
\href{https://math.berkeley.edu/~hutching/teach/proofs.pdf}{``Introduction
  to Mathematical Arguments''} \citep{Hutchings2003} અને
\href{https://eugeniacheng.com/wp-content/uploads/2017/02/cheng-proofguide.pdf}{``How to
  write proofs''} \citep{Cheng2004}.

\subsection{પ્રેરક વિડિયો}

ઘરકામ કરવા ઉત્સાહ નથી લાગતો? મન ઉદાસ છે? કદાચ આ વિડિયો
મદદરૂપ થાય!

\begin{itemize}
\item \url{https://www.youtube.com/watch?v=ZXsQAXx_ao0}
\item \url{https://www.youtube.com/watch?v=BQ4yd2W50No}
\item \url{https://www.youtube.com/watch?v=StTqXEQ2l-Y}
\end{itemize}
% END OLP-0612
\endgroup


\OLEndChapterHook
% END OLP-0602
\endgroup


\begingroup
\makeatletter\def\ol@sectioncs{section}\makeatother

% BEGIN OLP-0613
\olchapter{mth}{ind}{ગાણિતિક અનુમાન}
\olchapterid{mth}{ind}
\phantomsection\label{guunit:OLP-0613}


\begingroup
\makeatletter\def\ol@sectioncs{section}\makeatother

% BEGIN OLP-0614
\olfileid{mth}{ind}{int}

\olsection{પરિચય}
\phantomsection\label{guunit:OLP-0614}


આગમન સાબિતીની મહત્વની પદ્ધતિ છે. તે જુદાં જુદાં સ્વરૂપે
તર્કશાસ્ત્ર, સૈદ્ધાંતિક કમ્પ્યુટર વિજ્ઞાન અને ગણિતના લગભગ દરેક
ક્ષેત્રમાં વપરાય છે. તર્કશાસ્ત્રનાં ઘણાં પરિણામો સાબિત કરવા તેની
જરૂર પડે છે.

આગમનને ઘણી વાર નિગમનથી જુદું ગણવામાં આવે છે અને વિશેષ કિસ્સાઓમાંથી
સામાન્ય દાવા સુધી પહોંચતા અનુમાન તરીકે વર્ણવાય છે. દાખલા તરીકે,
આપણે ઘણા લીલા મરકત જોયા હોય અને જેને મરકત કહીએ એવું એક પણ
બિનલીલું ન જોયું હોય, તો બધા મરકત લીલા છે એવું તારવી શકીએ. આ
આગમનાત્મક તારણ છે: તે ઘણા વિશેષ કિસ્સાઓમાંથી (આ મરકત લીલું છે,
પેલું મરકત લીલું છે વગેરે) સામાન્ય દાવા (બધા મરકત લીલા છે) સુધી
પહોંચે છે. \emph{ગાણિતિક} આગમન પણ સામાન્ય દાવો સ્થાપે છે, પરંતુ
તે આ "સરળ આગમન"થી ઘણું જુદા પ્રકારનું છે.

ખૂબ સાદા શબ્દોમાં, ગણિતમાં આગમનથી કરેલી સાબિતી બતાવે છે કે
ચોક્કસ પ્રકારના બધા ગાણિતિક પદાર્થોમાં કોઈ એક ગુણધર્મ છે. સૌથી
સરળ કિસ્સામાં આ પદાર્થો પ્રાકૃતિક સંખ્યાઓ હોય છે. ત્યારે
આગમનથી કરેલી સાબિતી બધી પ્રાકૃતિક સંખ્યાઓમાં તે ગુણધર્મ છે
એવું સ્થાપે છે. તે આ રીતે કરે છે:
\begin{enumerate}
    \item $0$માં તે ગુણધર્મ છે, અને
    \item જ્યારે પણ કોઈ સંખ્યા~$k$માં તે ગુણધર્મ હોય, ત્યારે~$k+1$માં પણ હોય છે.
\end{enumerate}
પ્રાકૃતિક સંખ્યાઓ પરનું આગમન ઘણી વાર એવી ગાણિતિક વસ્તુઓ અંગેના
સામાન્ય દાવા સાબિત કરવા પણ વાપરી શકાય છે જેમને સંખ્યા આપી શકાય.
દાખલા તરીકે, દરેક સીમિત ગણમાં ઘટકોની સીમિત સંખ્યા~$n$ હોય છે.
જો આગમનથી બતાવી શકીએ કે દરેક સંખ્યા~$n$ માટે "કદ~$n$ના બધા
સીમિત ગણો \dots છે", તો બધા સીમિત ગણો વિશે કંઈક સાબિત થયું હશે.

આગમનને \emph{આગમનાત્મક રીતે વ્યાખ્યાયિત} ગાણિતિક પદાર્થો સુધી
પણ વિસ્તારી શકાય છે. દાખલા તરીકે, પ્રથમ-ક્રમ તર્ક જેવી
ઔપચારિક ભાષાની અભિવ્યક્તિઓ આગમનાત્મક રીતે વ્યાખ્યાયિત
થાય છે. \emph{સંરચનાત્મક આગમન} એવી બધી અભિવ્યક્તિઓ અંગે
પરિણામો સાબિત કરવાની રીત છે. ખાસ કરીને સંરચનાત્મક આગમન
તર્કશાસ્ત્રમાં ખૂબ ઉપયોગી છે અને વ્યાપકપણે વપરાય છે.
% END OLP-0614
\endgroup


\begingroup
\makeatletter\def\ol@sectioncs{section}\makeatother

% BEGIN OLP-0615
\olfileid{mth}{ind}{inN}

\olsection{$\Nat$ પર આગમન}
\phantomsection\label{guunit:OLP-0615}


સૌથી સરળ સ્વરૂપમાં આગમન બધી પ્રાકૃતિક સંખ્યાઓ માટેનાં પરિણામો
સાબિત કરવાની પદ્ધતિ છે. તેમાં એ હકીકત વપરાય છે કે $0$થી શરૂ
કરીને વારંવાર~$1$ ઉમેરતાં છેવટે દરેક પ્રાકૃતિક સંખ્યા મળે છે.
તેથી દરેક સંખ્યા માટે કંઈક સાચું છે તે સાબિત કરવા આપણે
(1)~તે $0$ માટે સાચું છે એમ સ્થાપી શકીએ અને (2)~જ્યારે તે
કોઈ સંખ્યા~$n$ માટે સાચું હોય ત્યારે પછીની સંખ્યા~$n+1$
માટે પણ સાચું છે એમ બતાવી શકીએ. "$n$માં ગુણધર્મ $P$ છે"
તેને ટૂંકમાં $P(n)$ લખીએ (અને "$k$માં ગુણધર્મ $P$ છે"ને
$P(k)$, વગેરે). તો દરેક $n \in \Nat$ માટે $P(n)$ છે એની
આગમનથી સાબિતીમાં આ બે ભાગ હોય છે:
\begin{enumerate}
\item $P(0)$ની સાબિતી, અને
\item કોઈ પણ~$k$ માટે, જો $P(k)$ તો $P(k+1)$ છે તેની સાબિતી.
\end{enumerate}
આ વાત સંપૂર્ણ સ્પષ્ટ થાય તે માટે ધારો કે (1) અને~(2) બંને
આપણી પાસે છે. ત્યારે (1) કહે છે કે $P(0)$ સાચું છે. (2) પણ
હોય તો ખાસ કરીને $P(0)$ હોય ત્યારે $P(0+1)$, એટલે $P(1)$,
થાય છે. આ "$k$ ગમે તે હોય, $P(k)$ હોય તો $P(k+1)$" એવા
સામાન્ય વિધાનમાં~$k$ માટે~$0$ મૂકવાથી મળે છે. તેથી
મોડસ પોનેન્સથી~$P(1)$ મળે છે. હવે ફરી (2)માં~$k$ માટે
$1$ મૂકતાં: $P(1)$ હોય તો~$P(2)$ મળે.\sourcecorrection{OLMD-163}{અહીં
  સ્થિર અંગ્રેજી મૂળમાં (2)ના $k$ને બદલે $n$ લખ્યું છે;
  (2)માં સર્વત્ર પરિમાણિત ચલ $k$ છે.} હમણાં જ~$P(1)$
સ્થાપ્યું હોવાથી મોડસ પોનેન્સથી~$P(2)$ મળે છે. આ રીતે
આગળ વધીએ. કોઈ પણ સંખ્યા~$n$ માટે આવું $n$~વાર કરતાં
છેવટે~$P(n)$ સુધી પહોંચીએ છીએ. આમ (1) અને~(2) મળીને
દરેક $n \in \Nat$ માટે~$P(n)$ સ્થાપે છે.

એક ઉદાહરણ જોઈએ. ધારો કે $n$ પાસા ફેંકતાં કેટલા જુદા
સરવાળા મળી શકે તે જાણવું છે. વિચિત્ર લાગે તો પણ $0$ પાસાથી
શરૂ કરીએ. એક પણ પાસો ન હોય તો ફેંકવાથી એક જ સરવાળો મળે:
એક પણ ટપકું નહીં, એટલે સરવાળો~$0$. તેથી શક્ય પરિણામોની
સંખ્યા~$1$ છે. ફક્ત એક પાસો હોય, એટલે $n=1$ હોય, તો
$1$થી~$6$ સુધીનાં છ મૂલ્યો મળે. બે પાસાથી $2$થી $12$
સુધીનો કોઈ પણ સરવાળો મળે, એટલે $11$ શક્યતાઓ. ત્રણ
પાસાથી $3$થી $18$ સુધીનાં મૂલ્યો, એટલે $16$ શક્યતાઓ
મળે. $1$, $6$, $11$, $16$: કોઈ નિયમ દેખાય છે; જવાબ
કદાચ $5n+1$ હશે? ખરેખર, $5n+1$થી વધારે મૂલ્યો શક્ય જ
નથી, કારણ કે $n$ પાસાથી મળતો લઘુતમ સરવાળો $n$ (બધા
પાસા પર $1$) અને મહત્તમ સરવાળો $6n$ (બધા પર $6$) છે,
અને તેમની વચ્ચે, છેડા સહિત, ફક્ત $5n+1$ સંખ્યાઓ છે.

\begin{thm}
  $n$ પાસાથી $n$ અને $6n$ વચ્ચેના બધા $5n+1$ શક્ય મૂલ્યો
  મેળવી શકાય છે.
\end{thm}

\begin{proof}
$P(n)$ એ "$n$ પાસા વડે $n$ અને~$6n$ વચ્ચેની કોઈ પણ
સંખ્યા મેળવી શકાય છે" એવો દાવો રાખીએ. આગમન વાપરવા
માટે આપણે સાબિત કરીએ:
\begin{enumerate}
\item \emph{આગમનનો આધાર} $P(1)$, એટલે ફક્ત એક પાસાથી
  $1$ અને $6$ વચ્ચેની દરેક સંખ્યા મેળવી શકાય છે.
\item \emph{આગમન-પગલું}: બધા $k$ માટે, જો $P(k)$ તો~$P(k+1)$.
\end{enumerate}

(1) $6$-બાજુવાળો પાસો જોઈને સાબિત થાય છે. તેને છ બાજુઓ છે
અને $1$થી~$6$ સુધીની દરેક સંખ્યા તેની કોઈ એક બાજુ પર છે.
તેથી એક જ પાસો વાપરીને $1$ અને~$6$ વચ્ચેની કોઈ પણ સંખ્યા
મેળવી શકાય છે.

(2) સાબિત કરવા આપણે શરતી વિધાનનું પૂર્વપદ, એટલે $P(k)$,
ધારીએ. આ ધારણાને \emph{આગમન-પરિકલ્પના} કહે છે. તેનો
ઉપયોગ કરીને~$P(k+1)$ સાબિત કરીએ. મુશ્કેલ ભાગ એ છે કે
$k+1$ પાસા ફેંકતાં મળતાં મૂલ્યોને $k$~પાસાથી મળતાં
મૂલ્યો અને વધારાના $k+1$મા પાસાથી મળતાં મૂલ્યોની
દૃષ્ટિએ કેવી રીતે વિચારવાં. આગમન-પરિકલ્પના વાપરવી
હોય તો આવું કરવું જ પડે.

આગમન-પરિકલ્પના મુજબ $k$~પાસાથી $k$ અને~$6k$ વચ્ચેની
કોઈ પણ સંખ્યા મેળવી શકાય છે. $(k+1)$મા પાસા પર~$1$
આવે તો કુલ સરવાળામાં $1$ ઉમેરાય. તેથી $k$~પાસા ફેંકી,
પછી $(k+1)$મા પાસા પર~$1$ મેળવીને $k+1$થી $6k+1$
સુધીનું કોઈ પણ મૂલ્ય મેળવી શકીએ. હવે શું બાકી છે?
$6k+2$થી $6k+6$ સુધીનાં મૂલ્યો. પહેલા $k$ પાસા પર
$6$ અને પછી $(k+1)$મા પાસા પર $2$થી $6$ વચ્ચેની
સંખ્યા મેળવીને આ મૂલ્યો મળે છે. આમ $k+1$ પાસાથી
$k+1$ અને $6(k+1)$ વચ્ચેની બધી સંખ્યાઓ મેળવી શકાય છે.
એટલે આગમન-પરિકલ્પના~$P(k)$ ધારીને~$P(k+1)$ સાબિત થયું.
\end{proof}

ઘણી વાર વસ્તુઓની (સંખ્યાઓ, ગણો વગેરેની) એવી શ્રેણી
વિશે કંઈક સાબિત કરવા આપણે આગમન વાપરીએ છીએ જેને
પોતાને "આગમનાત્મક રીતે" વ્યાખ્યાયિત કરી હોય, એટલે
$(n+1)$મી વસ્તુને $n$મી વસ્તુ વડે વ્યાખ્યાયિત કરી હોય.
દાખલા તરીકે,~$n$ સુધીની પ્રાકૃતિક સંખ્યાઓનો સરવાળો~$s_n$
આ રીતે વ્યાખ્યાયિત કરી શકીએ:
\begin{align*}
  s_0 & = 0\\
  s_{n+1} & = s_n + (n+1)
\end{align*}
આ વ્યાખ્યાથી મળે છે:
\begin{align*}
  s_0 & = 0,\\
  s_1 & = s_0 + 1 && = 1,\\
  s_2 & = s_1 + 2 && = 1 + 2 = 3\\
  s_3 & = s_2 + 3 && = 1 + 2 + 3 = 6, \text{ વગેરે}
\end{align*}
હવે આગમનથી સાબિત કરી શકીએ કે $s_n = n(n+1)/2$.

\begin{prop}
  $s_n = n(n+1)/2$.
\end{prop}

\begin{proof}
  આપણે (1)~$s_0 = 0\cdot(0 + 1)/2$ અને (2)~જો $s_k =
  k(k+1)/2$ તો $s_{k+1} = (k+1)(k+2)/2$ સાબિત કરવાનું છે.
  (1) સ્પષ્ટ છે. (2) માટે આગમન-પરિકલ્પના $s_k = k(k+1)/2$
  ધારીએ. તેને વાપરીને $s_{k+1} = (k+1)(k+2)/2$ બતાવવું છે.

  $s_{k+1}$ શું છે? વ્યાખ્યા મુજબ $s_{k+1} = s_k + (k+1)$.
  આગમન-પરિકલ્પના મુજબ $s_k = k(k+1)/2$. તેને અગાઉના
  સમીકરણમાં મૂકીએ તો માત્ર થોડું અપૂર્ણાંકોનું ગણિત બાકી રહે છે:
  \begin{align*}
    s_{k+1} & = \frac{k(k+1)}{2} + (k+1) = {}\\
    & = \frac{k(k+1)}{2} + \frac{2(k+1)}{2} = {}\\
    & = \frac{k(k+1) + 2(k+1)}{2} = {}\\
    & = \frac{(k+2)(k+1)}{2}.
  \end{align*}
\end{proof}

અહીં મુખ્ય પાઠ એ છે કે કોઈ આગમનાત્મક રીતે વ્યાખ્યાયિત
શ્રેણી~$a_n$ વિશે પરિણામ સાબિત કરવું હોય તો આગમન
સ્વાભાવિક પદ્ધતિ છે. અને જો શ્રેણી એવી રીતે વ્યાખ્યાયિત
ન પણ હોય (જેમ પાસા ફેંકવાના ઉદાહરણમાં), તોપણ~$k+1$નો
કિસ્સો~$k$ના કિસ્સા સાથે કોઈ રીતે જોડી શકાય તો આગમન
વાપરી શકાય છે.
% END OLP-0615
\endgroup


\begingroup
\makeatletter\def\ol@sectioncs{section}\makeatother

% BEGIN OLP-0616
\olfileid{mth}{ind}{str}

\olsection{પ્રબળ આગમન}
\phantomsection\label{guunit:OLP-0616}


ઉપર ચર્ચેલા આગમનના સિદ્ધાંતમાં આપણે $P(0)$ સાબિત કરીએ છીએ
અને જો $P(k)$ હોય તો $P(k+1)$ પણ હોય છે એ સાબિત કરીએ છીએ.
બીજા ભાગમાં $P(k)$ સાચું છે એમ ધારીને આ ધારણા વડે
$P(k+1)$ સાબિત કરીએ છીએ. અલબત્ત, તેની સમકક્ષ રીતમાં
$P(k-1)$ ધારીને $P(k)$ સાબિત કરી શકીએ. મુખ્ય વાત એ છે કે
કોઈ પણ સંખ્યાથી તેના અનુગામી સુધીનું તારણ કાઢી શકીએ:
કોઈ સંખ્યા માટે દાવો તેના પૂર્વગામી માટે સાચો છે એમ
ધારીને તે સંખ્યા માટે સાબિત કરી શકીએ.

આગમનના સિદ્ધાંતનો એક પ્રકાર એવો છે જેમાં $k$ના પૂર્વગામી
$k-1$ માટે જ દાવો સાચો છે એમ નથી ધારતા;~$k$થી નાની
બધી સંખ્યાઓ માટે તે સાચો છે એમ ધારીએ છીએ અને તે
ધારણાથી~$k$ માટે દાવો સ્થાપીએ છીએ. આથી પણ બધા~$n \in
\Nat$ માટે $P(n)$ મળે છે. કારણ કે એક વાર $P(0)$
સ્થાપીએ, ત્યારે~$1$થી નાની બધી સંખ્યાઓ માટે $P$ સાચો
છે એવું સ્થાપ્યું હોય છે. અને જો બધા $l<k$ માટે $P(l)$
હોય તો $P(k)$ થાય છે એમ જાણતા હોઈએ, તો ખાસ કરીને
$k=1$ માટે પણ આ જાણીએ છીએ. તેથી $P(1)$ તારવી
શકીએ. હવે $P(0)$ અને $P(1)$ સાબિત થયા, એટલે બધા
$l<2$ માટે $P(l)$ સાબિત થયું. બધા $l<2$ માટે $P(l)$
હોય તો $P(2)$ થાય છે એવું શરતી વિધાન પણ હોવાથી
$P(2)$ તારવી શકીએ, અને આ રીતે આગળ વધીએ.

ખરેખર, જો "બધા $k$ માટે, બધા $l<k$ માટે $P(l)$
હોય તો $P(k)$" એવું સામાન્ય શરતી વિધાન સ્થાપી શકીએ,
તો $P(0)$ અલગથી સ્થાપવાની જરૂર નથી, કારણ કે તે
શરતી વિધાનમાંથી જ નીકળે છે. યાદ રાખો કે "બધા
$l<k$ માટે $P(l)$" એવો સામાન્ય દાવો $l<k$ હોય
એવી કોઈ સંખ્યા ન હોય ત્યારે પણ સાચો છે. આ
ખાલી પરિમાણીકરણનો કિસ્સો છે: "$A$ પ્રકારની
બધી વસ્તુઓ $B$ પ્રકારની છે" એ જો $A$ પ્રકારની
કોઈ વસ્તુ ન હોય તો સાચું છે; કોઈ $x$ જો~$!A(x)$
સંતોષતું ન હોય, તો $\lforall[x][(!A(x) \lif !B(x))]$
સાચું છે. અહીં ઔપચારિક સ્વરૂપ
"$\lforall[l][(l < k \lif P(l))]$" થશે; અને જો
$l < k$ હોય એવી કોઈ સંખ્યા ન હોય તો તે સાચું છે.
હવે $k=0$ હોય ત્યારે બરાબર એવું જ થાય છે:
કોઈ $l<0$ નથી; તેથી "બધા $l<0$ માટે $P(l)$"
સાચું છે, $P$ ગમે તે હોય.\sourcecorrection{OLMD-164}{સ્થિર અંગ્રેજી મૂળમાં
  અહીં $P(l)$ને બદલે $P(0)$ લખાયું છે. ખાલી ક્ષેત્રે બંને
  સાર્વત્રિક વિધાનો સાચાં છે, પરંતુ આગમનના પૂર્વપદમાં
  દરેક $l<k$ માટે $P(l)$ જ જરૂરી છે.} તેથી "બધા
$l<k$ માટે $P(l)$ હોય તો $P(k)$" એવી સાબિતી
આપમેળે~$P(0)$ પણ સ્થાપે છે.

જો~$k$ માટે દાવો સ્થાપવા માત્ર $k-1$ માટેનો દાવો
પૂરતો ન હોય અને એક કે વધુ $l<k$ માટેનો દાવો
ધારવો પડે, તો આ પ્રકારનું આગમન ઉપયોગી છે.
% END OLP-0616
\endgroup


\begingroup
\makeatletter\def\ol@sectioncs{section}\makeatother

% BEGIN OLP-0617
\olfileid{mth}{ind}{idf}

\olsection{આગમનાત્મક વ્યાખ્યાઓ}
\phantomsection\label{guunit:OLP-0617}


તર્કશાસ્ત્રમાં આપણે ઘણી વાર પદાર્થોના પ્રકારોને
\emph{આગમનાત્મક રીતે} વ્યાખ્યાયિત કરીએ છીએ. એટલે
જે પ્રકારને વ્યાખ્યાયિત કરવો હોય તેમાં પદાર્થ ક્યારે
ગણાય તેના નિયમો આપીએ છીએ; એ નિયમો તે પ્રકારના
અગાઉના પદાર્થોમાંથી નવા પદાર્થો કેવી રીતે બનાવવા
તે પણ સમજાવે છે. દાખલા તરીકે, ભાષાનાં પદો અને
સૂત્રો જેવી ચિહ્નોની ખાસ જાતની શ્રેણીઓ આપણે
ઘણી વાર આગમનથી વ્યાખ્યાયિત કરીએ છીએ. એક સરળ
ઉદાહરણ તરીકે, $\mathrm{a}$, $\mathrm{b}$, $\mathrm{c}$,
$\mathrm{d}$ અક્ષરો,~$\circ$ ચિહ્ન અને $[$ તથા $]$
કૌંસથી બનેલી શ્રેણીઓ વિચારો; જેમ કે
"$[[\mathrm{c} \circ \mathrm{d}][$",
"$[\mathrm{a}[]\circ]$",
"$\mathrm{a}$" અથવા "$[[\mathrm{a} \circ \mathrm{b}]\circ
\mathrm{d}]$". પહેલી બે શ્રેણીઓમાં કંઈક "ખોટું"
લાગે: પહેલીમાં કૌંસોનો જરાય મેળ બેસતો નથી;
અને તમને એમ પણ લાગે કે "$\circ$"એ પોતે અર્થપૂર્ણ
હોય એવી અભિવ્યક્તિઓને "જોડવી" જોઈએ. ત્રીજી
અને ચોથી શ્રેણી વધુ યોગ્ય લાગે છે: દરેક "$[$"
માટે બંધ કરતો "$]$" છે (જો કૌંસો હોય તો), અને
દરેક $\circ$ની બંને બાજુ કૌંસોની જોડીથી ઘેરાયેલી
"યોગ્ય" અભિવ્યક્તિઓ મળે છે.

"સુઘડ પદ" ક્યારે ગણાય તે આપણે ચોક્કસ રીતે કહેવા
માંગીએ છીએ. પહેલું તો, દરેક અક્ષર પોતે સુઘડ છે.
ફક્ત એક અક્ષર ન હોય એવું કંઈ પણ "$[t \circ s]$"
સ્વરૂપનું હોવું જોઈએ, જ્યાં $s$ અને $t$ પોતે
સુઘડ હોય. ઊલટું, $t$ અને $s$ સુઘડ હોય તો તેમની
વચ્ચે $\circ$ મૂકીને અને આસપાસ કૌંસોની જોડી
મૂકીને નવું સુઘડ પદ બનાવી શકીએ. આવી ક્રિયાઓથી
સુઘડ પદોના ગણને \emph{વ્યાખ્યાયિત} કરી શકીએ.
આ \emph{આગમનાત્મક
  વ્યાખ્યા} છે.

\begin{defn}[સુઘડ પદો]
  \emph{સુઘડ પદો}નો ગણ નીચે મુજબ આગમનાત્મક રીતે
  વ્યાખ્યાયિત થાય છે:
  \begin{enumerate}
  \item કોઈ પણ અક્ષર $\mathrm{a}$, $\mathrm{b}$, $\mathrm{c}$,
    $\mathrm{d}$ સુઘડ પદ છે.
  \item $s_1$ અને $s_2$ સુઘડ પદો હોય, તો
    $[s_1 \circ s_2]$ પણ સુઘડ પદ છે.
  \item આ સિવાય બીજું કંઈ સુઘડ પદ નથી.
  \end{enumerate}
\end{defn}

આ વ્યાખ્યા કહે છે કે કોઈ વસ્તુ સુઘડ પદ ત્યારે અને
ત્યારે જ છે જ્યારે શરતો (1) અને~(2) મુજબ તેને
સીમિત સંખ્યાનાં પગલાંમાં બનાવી શકાય. પહેલા પગલામાં
ફક્ત એક-એક અક્ષરથી બનેલાં બધાં સુઘડ પદો બનાવીએ:
\[
\mathrm{a}, \mathrm{b}, \mathrm{c}, \mathrm{d}
\]
બીજા પગલામાં બનાવેલા પદો પર (2) લાગુ કરીએ. બે
અક્ષરોની બધી જોડીઓમાંથી મળે:
\[
[\mathrm{a} \circ \mathrm{a}], [\mathrm{a} \circ \mathrm{b}],
[\mathrm{b} \circ \mathrm{a}], \dots, [\mathrm{d} \circ \mathrm{d}]
\]
ત્રીજા પગલામાં અત્યાર સુધી બનાવેલાં કોઈ પણ બે
સુઘડ પદો પર ફરી (2) લાગુ કરીએ. તેથી
$[\mathrm{a} \circ [\mathrm{a} \circ
    \mathrm{a}]]$ જેવા નવા સુઘડ પદો મળે: અહીં $t$
પહેલા પગલાનું $\mathrm{a}$ છે અને $s$ બીજા પગલાનું
$[\mathrm{a} \circ \mathrm{a}]$ છે. વળી
$[[\mathrm{b} \circ
    \mathrm{c}] \circ [\mathrm{d} \circ \mathrm{b}]]$
એ બીજા પગલાનાં બે પદો $[\mathrm{b} \circ \mathrm{c}]$
અને $[\mathrm{d}
  \circ \mathrm{b}]$માંથી બને છે. આ રીતે આગળ
વધીએ. ખંડ (3) ખાતરી કરે છે કે આ રીતે ન બનેલું
કશું સુઘડ પદોના ગણમાં ઘૂસી ન જાય.

ધ્યાન આપો કે જે દરેક ચિહ્ન-શ્રેણી આપણને "સુઘડ"
લાગે છે તે આ વ્યાખ્યા મુજબ પણ સુઘડ છે એવું
આપણે હજી સાબિત કર્યું નથી. પરંતુ જે કંઈ આ
રીતે બનાવી શકીએ તે ખરેખર "સુઘડ લાગે" છે
એ સ્પષ્ટ હોવું જોઈએ: કૌંસોનો મેળ બેસે છે અને
$\circ$ પોતે સુઘડ હોય એવા ભાગોને જોડે છે.

આગમનાત્મક વ્યાખ્યાઓનું મુખ્ય લક્ષણ એ છે કે જો
બધાં સુઘડ પદો વિશે કંઈ સાબિત કરવું હોય, તો વ્યાખ્યા
કયા કિસ્સા વિચારવાના છે તે કહે છે. દાખલા તરીકે,
$t$ સુઘડ પદ છે એમ કહીએ ત્યારે આગમનાત્મક
વ્યાખ્યા તેનું શક્ય સ્વરૂપ જણાવે છે: $t$ અક્ષર
હોઈ શકે અથવા કોઈ સુઘડ પદો $s_1$ અને~$s_2$
માટે $t$ એ $[s_1 \circ s_2]$ હોઈ શકે. ખંડ (3)ને
લીધે આ જ બે શક્યતાઓ છે.

આગમનાત્મક રીતે વ્યાખ્યાયિત ગણની બધી વસ્તુઓ
વિશેના દાવા સાબિત કરતી વખતે આગમનનું પ્રબળ સ્વરૂપ
ખાસ મહત્વનું બને છે. દાખલા તરીકે, ધારો કે
લંબાઈ~$n$ના દરેક સુઘડ પદમાં $[$ની સંખ્યા~$< n/2$
છે એવું સાબિત કરવું છે. તેને બધા~$n$ વિશેના
દાવા તરીકે જોઈ શકાય છે: દરેક $n$ માટે લંબાઈ~$n$ના
કોઈ પણ સુઘડ પદમાં $[$ની સંખ્યા $< n/2$ છે.

\begin{prop}
  કોઈ પણ $n$ માટે લંબાઈ~$n$ના સુઘડ પદમાં
  $[$ની સંખ્યા $< n/2$ છે.
\end{prop}

\begin{proof}
આ પરિણામ (પ્રબળ) આગમનથી સાબિત કરવા નીચેનો
શરતી દાવો સાચો છે એમ બતાવવું પડે:
\begin{quote}
  જો દરેક $l < k$ માટે લંબાઈ~$l$ના કોઈ પણ
  સુઘડ પદમાં $[$ની સંખ્યા $< l/2$ હોય,
  તો લંબાઈ~$k$ના કોઈ પણ સુઘડ પદમાં
  $[$ની સંખ્યા $< k/2$ હોય.
\end{quote}
આ શરતી દાવો બતાવવા તેનું પૂર્વપદ સાચું ધારો:
કોઈ પણ $l<k$ માટે લંબાઈ~$l$ના સુઘડ પદોમાં
$[$ની સંખ્યા $< l/2$ છે. આ ધારણાને
આગમન-પરિકલ્પના કહીએ છીએ. લંબાઈ~$k$ના
સુઘડ પદો માટે પણ એ જ વાત બતાવવી છે.

ધારો કે $t$ એ લંબાઈ~$k$નું સુઘડ પદ છે.
સુઘડ પદો આગમનાત્મક રીતે વ્યાખ્યાયિત હોવાથી
બે કિસ્સા છે: (1)~$t$ પોતે એક અક્ષર છે,
અથવા (2)~કોઈ સુઘડ પદો $s_1$ અને~$s_2$
માટે $t$ એ $[s_1 \circ s_2]$ છે.
\begin{enumerate}
\item $t$ એક અક્ષર છે. ત્યારે $k = 1$ અને
$t$માં $[$ની સંખ્યા~$0$ છે. $0 < 1/2$
હોવાથી દાવો સાચો છે.
\item કોઈ સુઘડ પદો $s_1$ અને~$s_2$ માટે
$t$ એ $[s_1 \circ s_2]$ છે. $l_1$ એ~$s_1$ની
અને $l_2$ એ~$s_2$ની લંબાઈ માનો. ત્યારે $t$ની
લંબાઈ~$k$ એ $l_1+l_2+3$ છે ($s_1$ અને
$s_2$ની લંબાઈ ઉપરાંત $[$, $\circ$, $]$
એવાં ત્રણ ચિહ્નો). $l_1+l_2+3$ હંમેશાં
$l_1$થી મોટું હોવાથી $l_1 < k$. એ જ રીતે
$l_2 < k$. તેથી આગમન-પરિકલ્પના $s_1$ અને~$s_2$
બંને પદો પર લાગુ પડે છે: $s_1$માં $[$ની
સંખ્યા~$m_1$ એ $< l_1/2$ છે અને~$s_2$માં
$[$ની સંખ્યા~$m_2$ એ $< l_2/2$ છે.

  $t$માં $[$ની સંખ્યા એ~$s_1$માંની $[$ની
  સંખ્યા, વત્તા~$s_2$માંની $[$ની સંખ્યા,
  વત્તા~$1$ છે; એટલે $m_1 + m_2 + 1$.
  $m_1 < l_1/2$ અને $m_2 < l_2/2$
  હોવાથી મળે:
  \[
  m_1 + m_2 + 1 < \frac{l_1}{2} + \frac{l_2}{2} + 1 = \frac{l_1+l_2+2}{2} < \frac{l_1+l_2+3}{2} = k/2.
  \]
\end{enumerate}
દરેક કિસ્સામાં આપણે (આગમન-પરિકલ્પનાના આધારે)
બતાવ્યું કે $t$માં $[$ની સંખ્યા $< k/2$ છે.
પ્રબળ આગમનથી પ્રતિજ્ઞાપન ફલિત થાય છે.
\end{proof}

\begin{prob}
  અતિસુઘડ પદોનો ગણ આ રીતે વ્યાખ્યાયિત કરો:
  \begin{enumerate}
  \item કોઈ પણ અક્ષર $\mathrm{a}$, $\mathrm{b}$, $\mathrm{c}$,
    $\mathrm{d}$ અતિસુઘડ પદ છે.
  \item $s$ અતિસુઘડ પદ હોય તો $[s]$ પણ છે.
  \item $s_1$ અને $s_2$ અતિસુઘડ પદો હોય, તો
    $[s_1 \circ s_2]$ પણ છે.
  \item આ સિવાય બીજું કંઈ અતિસુઘડ પદ નથી.
  \end{enumerate}
  લંબાઈ~$n$ના અતિસુઘડ પદ~$t$માં $[$ની
  સંખ્યા $\le n/2 +1$ છે એમ બતાવો.
\end{prob}
% END OLP-0617
\endgroup


\begingroup
\makeatletter\def\ol@sectioncs{section}\makeatother

% BEGIN OLP-0618
\olfileid{mth}{ind}{sti}

\olsection{સંરચનાત્મક આગમન}
\phantomsection\label{guunit:OLP-0618}


અત્યાર સુધી આપણે બધી પ્રાકૃતિક સંખ્યાઓ વિશેનાં
પરિણામો સ્થાપવા આગમન વાપર્યું છે. પરંતુ અનુરૂપ
સિદ્ધાંત આગમનાત્મક રીતે વ્યાખ્યાયિત ગણના બધા
ઘટકો વિશે પરિણામો સાબિત કરવા સીધો વાપરી
શકાય છે. તેને ઘણી વાર \emph{સંરચનાત્મક} આગમન
કહે છે, કારણ કે તે આગમનાત્મક રીતે વ્યાખ્યાયિત
પદાર્થોની સંરચના પર આધાર રાખે છે.

સામાન્ય રીતે આગમનાત્મક વ્યાખ્યામાં (a)~ગણના
"શરૂઆતના" ઘટકોની યાદી અને (b)~ગણના
અગાઉના ઘટકોમાંથી નવા ઘટકો
બનાવતી ક્રિયાઓની યાદી હોય છે. દાખલા તરીકે,
સુઘડ પદોના કિસ્સામાં શરૂઆતના પદાર્થો અક્ષરો છે.
એક જ ક્રિયા છે:
\begin{align*}
  o(s_1, s_2) = & [s_1 \circ s_2]
\end{align*}
પ્રાકૃતિક સંખ્યાઓ~$\Nat$ને પણ આગમનાત્મક
વ્યાખ્યાથી મળેલી માની શકો: શરૂઆતનો પદાર્થ~$0$
છે અને ક્રિયા અનુગામી વિધેય~$x + 1$ છે.

આગમનાત્મક રીતે વ્યાખ્યાયિત ગણના બધા ઘટકો
વિશે કંઈક સાબિત કરવા, એટલે ગણના દરેક
ઘટકમાં ગુણધર્મ~$P$ છે એમ બતાવવા,
આપણે:
\begin{enumerate}
\item શરૂઆતના પદાર્થોમાં~$P$ છે એમ સાબિત કરવું, અને
\item દરેક ક્રિયા~$o$ માટે, તેના નિવેશોમાં~$P$
  હોય તો પરિણામ માટે પણ એવું જ છે એમ સાબિત કરવું.
\end{enumerate}
દાખલા તરીકે, બધાં સુઘડ પદો વિશે કંઈક સાબિત
કરવા પહેલાં બધા અક્ષરો માટે તે સાચું છે એમ
અને પછી $s_1$ અને $s_2$ બંને માટે તે સાચું
હોય તો $[s_1 \circ s_2]$ માટે પણ સાચું છે
એમ સાબિત કરીએ.

\begin{prop}
  કોઈ પણ સુઘડ પદ~$t$માં $[$ની સંખ્યા $]$ની
  સંખ્યા જેટલી જ છે.
\end{prop}

\begin{proof}
સંરચનાત્મક આગમન વાપરીએ. સુઘડ પદો આગમનાત્મક
રીતે વ્યાખ્યાયિત છે: અક્ષરો શરૂઆતના પદાર્થો છે
અને $o$ ક્રિયા જૂનાં સુઘડ પદોમાંથી નવાં બનાવે છે.
\begin{enumerate}
\item દરેક અક્ષર માટે દાવો સાચો છે, કારણ કે એક
  અક્ષરમાં $[$ની સંખ્યા~$0$ છે અને $]$ની
  સંખ્યા પણ~$0$ છે.
\item ધારો કે $s_1$માં $[$ની સંખ્યા $]$ની
  સંખ્યા જેટલી છે, અને $s_2$ માટે પણ એવું
  જ છે. $o(s_1, s_2)$માં, એટલે
  $[s_1 \circ s_2]$માં, $[$ની સંખ્યા એ
  $s_1$ અને $s_2$માંના $[$ની સંખ્યાનો
  સરવાળો વત્તા એક છે. $o(s_1, s_2)$માં
  $]$ની સંખ્યા એ $s_1$ અને $s_2$માંના
  $]$ની સંખ્યાનો સરવાળો વત્તા એક છે.
  તેથી $o(s_1, s_2)$માં $[$ની સંખ્યા
  $o(s_1,s_2)$માં $]$ની સંખ્યા જેટલી છે.
\end{enumerate}
\end{proof}

\begin{prob}
  કોઈ પણ સુઘડ પદની શરૂઆત~$]$થી થતી નથી
  એ સંરચનાત્મક આગમનથી સાબિત કરો.
\end{prob}

સંરચનાત્મક આગમનથી બીજી સાબિતી જોઈએ:
ચિહ્ન-શ્રેણી~$t$નો અરિક્ત ઉચિત આરંભિક ખંડ
એવી શ્રેણી~$s$ છે જે ડાબેથી વાંચતાં ચિહ્ને
ચિહ્ને~$t$ સાથે મળે છે, પરંતુ $t$ વધુ લાંબી
છે. દાખલા તરીકે, $[a \circ {}$ એ
$[a \circ b]$નો ઉચિત આરંભિક ખંડ છે; પરંતુ
$[b \circ {}$ નથી (બીજા ચિહ્ને ભેદ છે)
અને $[a \circ b]$ પણ નથી (લંબાઈ સરખી છે).
\sourcecorrection{OLMD-165}{સ્થિર અંગ્રેજી મૂળમાં ઉચિત આરંભિક
  ખંડ અરિક્ત હોવાની શરત કહેવાઈ નથી. ખાલી શ્રેણી
  પણ ઉચિત આરંભિક ખંડ ગણીએ તો આગળનું પ્રતિજ્ઞાપન
  ખોટું થાય: તેમાં $[$ અને $]$ બંનેની સંખ્યા શૂન્ય છે.
  આગળની સાબિતીમાં ખાલી ખંડનો કિસ્સો નથી; અહીં
  અભિપ્રેત અરિક્ત શરત સ્પષ્ટ કરી છે.}

\begin{prop}\ollabel{prop:initial}
  સુઘડ પદ~$t$ના દરેક અરિક્ત ઉચિત આરંભિક ખંડમાં
  $[$ની સંખ્યા $]$ની સંખ્યાથી વધુ છે.
\end{prop}

\begin{proof}
  $t$ પર આગમન કરીએ:
  \begin{enumerate}
  \item $t$ પોતે એક અક્ષર છે: ત્યારે $t$નો
    કોઈ અરિક્ત ઉચિત આરંભિક ખંડ નથી.
  \item કોઈ સુઘડ પદો $s_1$ અને~$s_2$ માટે
    $t = [s_1 \circ s_2]$ છે. $r$ જો $t$નો
    અરિક્ત ઉચિત આરંભિક ખંડ હોય, તો નીચેની
    શક્યતાઓ છે:
    \begin{enumerate}
    \item $r$ ફક્ત $[$ છે: ત્યારે $r$માં
      $]$ કરતાં એક $[$ વધુ છે.
    \item $r$ એ $[r_1$ છે, જ્યાં $r_1$ એ~$s_1$નો
      અરિક્ત ઉચિત આરંભિક ખંડ છે: $s_1$ સુઘડ પદ
      હોવાથી આગમન-પરિકલ્પના મુજબ $r_1$માં
      $]$ કરતાં વધુ $[$ છે અને $[r_1$માં
      પણ એવું જ છે.
    \item $r$ એ $[s_1$ અથવા $[s_1 \circ {}$
      છે: અગાઉના પરિણામ મુજબ~$s_1$માં $[$
      અને $]$ની સંખ્યા સરખી છે; તેથી
      $[s_1$ અથવા $[s_1 \circ {}$માં
      $]$ કરતાં એક $[$ વધુ છે.
    \item $r$ એ $[s_1 \circ r_2$ છે, જ્યાં
      $r_2$ એ~$s_2$નો અરિક્ત ઉચિત આરંભિક ખંડ છે:
      આગમન-પરિકલ્પના મુજબ $r_2$માં $]$
      કરતાં વધુ $[$ છે. અગાઉના પરિણામ
      મુજબ~$s_1$માં $[$ અને $]$ની સંખ્યા
      સરખી છે. તેથી $[s_1 \circ r_2$માં
      $]$ કરતાં વધુ $[$ છે.
    \item $r$ એ $[s_1 \circ s_2$ છે:
      અગાઉના પરિણામ મુજબ $s_1$માં $[$
      અને $]$ની સંખ્યા સરખી છે, અને~$s_2$
      માટે પણ એવું જ છે. તેથી
      $[s_1 \circ s_2$માં $]$ કરતાં એક
      $[$ વધુ છે.
    \end{enumerate}
  \end{enumerate}
\end{proof}
% END OLP-0618
\endgroup


\begingroup
\makeatletter\def\ol@sectioncs{section}\makeatother

% BEGIN OLP-0619
\olfileid{mth}{ind}{rel}

\olsection{સંબંધો અને વિધેયો}
\phantomsection\label{guunit:OLP-0619}


પદાર્થોનો ગણ (જેમ પ્રાકૃતિક સંખ્યાઓ અથવા સુઘડ પદો)
આગમનાત્મક રીતે વ્યાખ્યાયિત કર્યો હોય ત્યારે આ પદાર્થો
\emph{પરના સંબંધો} પણ આગમનથી વ્યાખ્યાયિત કરી શકીએ.
દાખલા તરીકે, આ વિચાર લો: સુઘડ પદ~$t_1$ જો સુઘડ
પદ~$t_2$ના ભાગરૂપે આવતું હોય તો તે તેનું ઉપપદ છે.
તે માટે $t_1 \sqsubseteq t_2$ ચિહ્ન વાપરીએ.
અલબત્ત, દરેક સુઘડ પદ પોતાનું ઉપપદ છે:
$t \sqsubseteq t$. આ સંબંધની આગમનાત્મક વ્યાખ્યા
આ રીતે આપી શકીએ:

\begin{defn}
  સુઘડ પદ~$t_1$ એ~$t_2$નું ઉપપદ છે એવો સંબંધ,
  $t_1 \sqsubseteq t_2$,~$t_2$ પરના આગમનથી
  નીચે મુજબ વ્યાખ્યાયિત થાય છે:
  \begin{enumerate}
  \item $t_2$ અક્ષર હોય તો $t_1 \sqsubseteq t_2$
    ત્યારે અને ત્યારે જ જ્યારે $t_1 = t_2$.
  \item $t_2$ એ $[s_1 \circ s_2]$ હોય તો
    $t_1 \sqsubseteq t_2$ ત્યારે અને ત્યારે જ જ્યારે
    $t_1 = t_2$, $t_1 \sqsubseteq s_1$ અથવા
    $t_1 \sqsubseteq s_2$.
  \end{enumerate}
\end{defn}

આ વ્યાખ્યાથી, દાખલા તરીકે, $\mathrm{a}
\sqsubseteq [\mathrm{b} \circ \mathrm{a}]$
મળે છે. કારણ કે (2) કહે છે કે
$\mathrm{a} \sqsubseteq [\mathrm{b} \circ \mathrm{a}]$
ત્યારે અને ત્યારે જ જ્યારે
$\mathrm{a} = [\mathrm{b} \circ \mathrm{a}]$,
અથવા $\mathrm{a} \sqsubseteq \mathrm{b}$,
અથવા $\mathrm{a} \sqsubseteq \mathrm{a}$.
પહેલી બે વાતો ખોટી છે: $\mathrm{a}$ એ
$[\mathrm{b} \circ
  \mathrm{a}]$ નથી, અને (1) મુજબ
$\mathrm{a} \sqsubseteq \mathrm{b}$
ત્યારે અને ત્યારે જ જ્યારે $\mathrm{a} = \mathrm{b}$,
જે પણ ખોટું છે. પરંતુ (1) મુજબ જ
$\mathrm{a} \sqsubseteq \mathrm{a}$
ત્યારે અને ત્યારે જ જ્યારે $\mathrm{a} = \mathrm{a}$,
જે સાચું છે.

ધ્યાન આપવાની અગત્યની વાત એ છે કે આ વ્યાખ્યા
સફળ થવી એક એવી હકીકત પર આધાર રાખે છે જે
આપણે હજી સાબિત કરી નથી: દરેક સુઘડ પદ~$t$
કાં તો પોતે એક અક્ષર હોય છે, અથવા
\emph{એકમાત્ર રીતે નિર્ધારિત} સુઘડ પદો
$s_1$ અને $s_2$ મળે છે જેથી
$t = [s_1 \circ s_2]$. અહીં "એકમાત્ર રીતે
નિર્ધારિત"નો અર્થ એ છે કે
$t = [s_1 \circ s_2]$ હોય તો
$s_1 \neq r_1$ અથવા $s_2 \neq r_2$ હોય
એવા બીજા બે પદો માટે તે \emph{સાથે સાથે}
$= [r_1 \circ r_2]$ પણ નહીં હોય. જો એવું
હોય તો ખંડ~(2) પોતાનાં જ પરિણામો સાથે
વિરોધમાં આવી શકે:~$t_2$ને
$[s_1 \circ s_2]$ તરીકે વાંચીએ ત્યારે
$t_1 \sqsubseteq t_2$ મળે, પણ~$t_2$ને
$[r_1 \circ r_2]$ તરીકે વાંચીએ ત્યારે
$t_1 \sqsubseteq t_2$ ન મળે. આવું ન
થઈ શકે તે સાબિત કરતાં પહેલાં જ્યાં
આવું \emph{થઈ શકે} તેવું ઉદાહરણ જોઈએ.

\begin{defn}
  \emph{કૌંસવિહોણાં પદો} આગમનાત્મક રીતે
  આ પ્રમાણે વ્યાખ્યાયિત કરો:
  \begin{enumerate}
  \item દરેક અક્ષર કૌંસવિહોણું પદ છે.
    \item $s_1$ અને $s_2$ કૌંસવિહોણાં પદો
      હોય તો $s_1 \circ s_2$ પણ કૌંસવિહોણું
      પદ છે.
    \item આ સિવાય બીજું કંઈ કૌંસવિહોણું પદ નથી.
  \end{enumerate}
\end{defn}

કૌંસવિહોણાં પદોનાં ઉદાહરણો
$\mathrm{a}$, $\mathrm{b} \circ
\mathrm{d}$ અને $\mathrm{b} \circ \mathrm{a}
\circ \mathrm{b}$ છે. હવે કૌંસવિહોણાં
પદો માટે પણ ઉપરની જેમ "ઉપપદ" વ્યાખ્યાયિત
કરીએ તો બીજો ખંડ આવો થશે:
\begin{center}
  જો $t_2 = s_1 \circ s_2$ હોય, તો
  $t_1 \sqsubseteq t_2$ ત્યારે અને ત્યારે જ
  જ્યારે $t_1 = t_2$, $t_1
  \sqsubseteq s_1$ અથવા $t_1 \sqsubseteq s_2$.
\end{center}

હવે $\mathrm{b} \circ \mathrm{a} \circ \mathrm{b}$
એ $s_1 \circ s_2$ સ્વરૂપનું છે, જ્યાં
\begin{align*}
  s_1 & = \mathrm{b} \text{ અને} & s_2 & = \mathrm{a} \circ \mathrm{b}.
\intertext{તે $r_1 \circ r_2$ સ્વરૂપનું પણ છે, જ્યાં}
  r_1 & = \mathrm{b} \circ \mathrm{a} \text{ અને} & r_2 &= \mathrm{b}.
\end{align*}
હવે $\mathrm{a} \circ \mathrm{b}$ એ
$\mathrm{b} \circ
\mathrm{a} \circ \mathrm{b}$નું ઉપપદ છે?
પહેલા વાંચન પ્રમાણે જવાબ હા છે, અને બીજા
વાંચન પ્રમાણે ના છે.

સુઘડ પદને બીજાં સુઘડ પદોમાંથી બનાવવાની રીત
એકમાત્ર હોય તે ગુણધર્મને \emph{એકમાત્ર વાંચન}
કહે છે. આવા આગમનાત્મક રીતે વ્યાખ્યાયિત
પદાર્થો પર સંબંધોની આગમનાત્મક વ્યાખ્યાઓ
મહત્વની હોવાથી આ ગુણધર્મ સાબિત કરવો પડે.

\begin{prop}
  ધારો કે $t$ સુઘડ પદ છે. ત્યારે કાં તો $t$
  પોતે એક અક્ષર છે, અથવા એકમાત્ર રીતે
  નિર્ધારિત સુઘડ પદો $s_1$, $s_2$ મળે છે
  જેથી~$t = [s_1 \circ s_2]$.
\end{prop}

\begin{proof}
  $t$ પોતે અક્ષર હોય તો શરત પૂરી થાય છે.
  તેથી ધારો કે $t$ પોતે અક્ષર નથી.
  આગમનાત્મક વ્યાખ્યા મુજબ ત્યારે કોઈ
  સુઘડ પદો $s_1$ અને~$s_2$ માટે $t$
  એ $[s_1 \circ s_2]$ સ્વરૂપનું હોવું જોઈએ.
  હવે આ બંને એકમાત્ર રીતે નિર્ધારિત છે
  તે બતાવવું બાકી છે: જો
  $t = [r_1 \circ r_2]$ હોય તો
  $s_1 = r_1$ અને $s_2 = r_2$.

  તેથી ધારો કે સુઘડ પદો $s_1$, $s_2$,
  $r_1$, $r_2$ માટે
  $t = [s_1 \circ s_2]$ અને
  $t = [r_1 \circ r_2]$ બંને છે.
  $s_1 = r_1$ અને $s_2 = r_2$ બતાવવું
  છે. પહેલા $s_1$ અને $r_1$ એકસમાન
  હોવા જોઈએ, કારણ કે નહિ તો એક બીજાનો
  ઉચિત આરંભિક ખંડ છે. પરંતુ
  \olref[sti]{prop:initial} મુજબ $s_1$
  અને~$r_1$ બંને સુઘડ પદો હોય ત્યારે
  એવું અશક્ય છે.
  હવે $s_1 = r_1$ હોય તો સ્પષ્ટપણે
  $s_2 = r_2$ પણ છે.
\end{proof}

વિધેયો પણ આગમનાત્મક રીતે વ્યાખ્યાયિત કરી
શકીએ. દાખલા તરીકે, કોઈ પણ સુઘડ પદમાં
એકની અંદર એક આવેલા $[\dots]$ કૌંસોનું
મહત્તમ ઊંડાણ આપતું વિધેય~$f$ આ રીતે
વ્યાખ્યાયિત કરીએ:

\begin{defn}
  \ollabel{defn:depth} સુઘડ પદનું
  \emph{ઊંડાણ}, $f(t)$, આગમનાત્મક રીતે
  નીચે મુજબ વ્યાખ્યાયિત થાય છે:
  \[
  f(t) = \begin{cases}
    0 & \text{ જો $t$ અક્ષર હોય}\\
    \max(f(s_1), f(s_2)) + 1 & \text{ જો $t = [s_1 \circ s_2]$ હોય.}
  \end{cases}
  \]
\end{defn}

દાખલા તરીકે,
\begin{align*}
  f([\mathrm{a} \circ \mathrm{b}]) & =
    \max(f(\mathrm{a}),f(\mathrm{b})) + 1 = \\
  &= \max(0, 0) + 1 = 1, \text{ અને}\\
  f([[\mathrm{a} \circ \mathrm{b}] \circ \mathrm{c}]) & =
    \max(f([\mathrm{a} \circ \mathrm{b}]), f(\mathrm{c})) + 1 = \\
  & = \max(1,0) + 1 = 2.
\end{align*}

અહીં આપણે અલબત્ત $s_1$ અને $s_2$ સુઘડ
પદો છે એમ ધારીએ છીએ અને દરેક સુઘડ પદ
કાં તો અક્ષર છે અથવા $[s_1 \circ s_2]$
સ્વરૂપનું છે તે હકીકત વાપરીએ છીએ. આ
સ્વરૂપની રચના એક જ રીતે થઈ શકે એ
અહીં પણ અગત્યનું છે. તેનું કારણ જોવા
અગાઉ વ્યાખ્યાયિત કૌંસવિહોણાં પદો ફરી
વિચારો. તેમને માટે અનુરૂપ "વ્યાખ્યા"
આવી થશે:
\[
  g(t) =
  \begin{cases}
    0 & \text{ જો $t$ અક્ષર હોય}\\
   \max(g(s_1), g(s_2)) + 1 & \text{ જો $t = s_1 \circ s_2$ હોય.}
  \end{cases}
\]
હવે કૌંસવિહોણું પદ
$\mathrm{a} \circ \mathrm{b} \circ
\mathrm{c} \circ \mathrm{d}$ વિચારો.
તેને એક કરતાં વધુ રીતે વાંચી શકાય છે;
દાખલા તરીકે, $s_1 \circ s_2$ સ્વરૂપમાં
\begin{align*}
  s_1 & = \mathrm{a} \text{ અને} &
  s_2 & = \mathrm{b} \circ \mathrm{c} \circ \mathrm{d},
\intertext{અથવા $r_1 \circ r_2$ સ્વરૂપમાં}
  r_1 & = \mathrm{a} \circ \mathrm{b} \text{ અને} &
  r_2 &= \mathrm{c} \circ \mathrm{d}.
\end{align*}
પહેલા વાંચન મુજબ $g$ ગણીએ તો મળે:
\begin{align*}
  g(s_1 \circ s_2) & =
  \max(g(\mathrm{a}), g(\mathrm{b} \circ \mathrm{c} \circ \mathrm{d})) + 1 =\\ & =
  \max(0,2) + 1 = 3
  \intertext{જ્યારે બીજા વાંચન મુજબ મળે}
  g(r_1 \circ r_2) & =
  \max(g(\mathrm{a} \circ \mathrm{b}), g(\mathrm{c} \circ \mathrm{d})) + 1=\\ &
  = \max(1,1) + 1 = 2
\end{align*}
પરંતુ વિધેય હંમેશાં એકમાત્ર મૂલ્ય
આપતું હોવું જોઈએ; તેથી~$g$ની આપણી
"વ્યાખ્યા" ખરેખર વિધેય વ્યાખ્યાયિત જ કરતી નથી.
\sourcecorrection{OLMD-166}{સ્થિર અંગ્રેજી મૂળમાં બે
  જગ્યાએ અક્ષર~$\mathrm{b}$ને ચલ~$b$ની
  જેમ ટાઇપસેટ કર્યું છે: ઉપપદના ઉદાહરણમાં
  અને બીજા કૌંસવિહોણા વાંચનના $r_1$માં.
  મૂળનો અભિપ્રાય અક્ષર~$\mathrm{b}$ જ છે;
  બંને જગ્યાએ તે ચિહ્ન સ્પષ્ટ કર્યું છે.}

\begin{prob}
  વિધેય $l$ની આગમનાત્મક વ્યાખ્યા આપો, જ્યાં
  $l(t)$ એ સુઘડ પદ~$t$માં ચિહ્નોની સંખ્યા છે.
\end{prob}

\begin{prob}
  સુઘડ પદો $t$ પર સંરચનાત્મક આગમનથી
  $f(t) < l(t)$ સાબિત કરો (અહીં $l(t)$
  એ~$t$માં ચિહ્નોની સંખ્યા છે અને $f(t)$
  એ \olref[mth][ind][rel]{defn:depth}માં
  વ્યાખ્યાયિત~$t$નું ઊંડાણ છે).
\end{prob}
% END OLP-0619
\endgroup


\OLEndChapterHook
% END OLP-0613
\endgroup


\OLEndPartHook
% END OLP-0601
